% !TeX spellcheck = pt_BR
\documentclass[12pt,portuguese]{article}

%----------------------------------------------------------------------
% PACOTES ESSENCIAIS E TRADUÇÃO
%----------------------------------------------------------------------
\usepackage[brazilian]{babel}     % Suporte à língua portuguesa
\usepackage[utf8]{inputenc}       % Codificação de entrada
\usepackage[T1]{fontenc}          % Codificação da fonte
\usepackage{lmodern}              % Fonte escalável

%----------------------------------------------------------------------
% PACOTES MATEMÁTICOS
%----------------------------------------------------------------------
\usepackage{amsmath, amsthm, amssymb, amstext, mathtools}
\usepackage{mathrsfs}             % Para fontes como \mathscr
\usepackage{dsfont}               % Para o indicador de conjunto \mathds{1}
\usepackage[all]{xy}              % Para diagramas comutativos
\usepackage{imakeidx}             % Criação de Lista de Símbolos / Índice

% Desigualdades sempre com símbolos inclinados
\let\leq\leqslant
\let\le\leqslant
\let\ge\geqslant
\let\geq\geqslant

%----------------------------------------------------------------------
% FORMATAÇÃO, LAYOUT E CORES
%----------------------------------------------------------------------
\usepackage{geometry}
\usepackage{varwidth}
\usepackage{graphicx}
\usepackage[pdftex]{xcolor}
\usepackage{titlesec}
\usepackage{appendix}
\usepackage{enumitem}
\usepackage{float}
\usepackage{booktabs}
\usepackage{array}
\usepackage{colortbl}
\usepackage{fontawesome5}
\usepackage{lucide-icons}
\usepackage{caption}
\usepackage{stackengine}
\captionsetup{width=0.8\textwidth}

% Definição de Cores
\definecolor{teal}{RGB}{15,118,110}
\definecolor{text}{RGB}{28,25,23}
\definecolor{muted}{RGB}{87,83,78}
\definecolor{card}{RGB}{255,255,255}
\definecolor{border}{RGB}{231,229,228}

% Definição de margens e espaçamento
\setlength{\topmargin}{-0.7in}
\setlength{\oddsidemargin}{0in}
\setlength{\textheight}{9.55in}
\setlength{\textwidth}{6.5in}
\setlength{\baselineskip}{12mm}

% Formatação dos títulos de seção
\titleformat{\section}{\Large\bfseries}{\thesection.}{0.5em}{}
\titleformat{\subsection}{\large\bfseries}{\thesubsection.}{0.5em}{}
\titleformat{\subsubsection}{\normalfont\itshape}{\thesubsubsection.}{0.6em}{}

%----------------------------------------------------------------------
% TIKZ E GRÁFICOS
%----------------------------------------------------------------------
\usepackage{tikz}
\usepackage{pgfplots}
\pgfplotsset{compat=1.17}
\usetikzlibrary{positioning,calc,fit,nfold,decorations.pathreplacing,calligraphy,tikzmark,patterns,arrows.meta,bending,hobby}

%----------------------------------------------------------------------
% REFERÊNCIAS, LINKS E METADADOS DO PDF
%----------------------------------------------------------------------
\usepackage[pdftex, colorlinks, linkcolor={red!50!black}, 
citecolor={red}, urlcolor={blue!80!black}]{hyperref}

\hypersetup{
pdftitle={Desigualdades de Markov e Chebyshev e o lema de Borel-Cantelli},
pdfauthor={L. Cioletti},
pdfsubject={Desigualdades de Markov, Chebyshev, Jensen e Lyapunov, convexidade e lema de Borel-Cantelli},
pdfkeywords={desigualdade de Markov, desigualdade de Chebyshev, convexidade, desigualdade de Jensen, desigualdade de Lyapunov, lema de Borel-Cantelli, MAT-UnB}
}

% --- PACOTE E COMANDO PARA A LICENÇA CREATIVE COMMONS ---
\usepackage[type={CC},modifier={by-nc-sa},version={4.0}]{doclicense}
\newcommand\blfootnote[1]{%
\begingroup
\renewcommand\thefootnote{}\footnote{#1}%
\addtocounter{footnote}{-1}%
\endgroup
}

%----------------------------------------------------------------------
% ALIASCNT E AMBIENTES DE TEOREMA (Cleveref)
%----------------------------------------------------------------------
\usepackage{aliascnt}
\theoremstyle{definition}
\newtheorem{theorem}{Teorema}

\newaliascnt{definition}{theorem}
\newtheorem{definition}[definition]{Definição}
\aliascntresetthe{definition}

\newaliascnt{example}{theorem}
\newtheorem{example}[example]{Exemplo}
\aliascntresetthe{example}

\newaliascnt{proposition}{theorem}
\newtheorem{proposition}[proposition]{Proposição}
\aliascntresetthe{proposition}

\newaliascnt{corollary}{theorem}
\newtheorem{corollary}[corollary]{Corolário}
\aliascntresetthe{corollary}

\newaliascnt{lemma}{theorem}
\newtheorem{lemma}[lemma]{Lema}
\aliascntresetthe{lemma}

\newaliascnt{remark}{theorem}
\newtheorem{remark}[remark]{Observação}
\aliascntresetthe{remark}

\newtheoremstyle{plainexercise}
  {3pt}{3pt}{\normalfont}{}{\bfseries}{.}{.5em}
  {\thmnumber{#2}\thmnote{ #3}}
\theoremstyle{plainexercise}
\newtheorem{exercise}{Exercício}

% Versão marcada dos exercícios: a estrela fica pendurada na margem esquerda,
% como nas notas anteriores.
\newtheoremstyle{starredexercise}
  {3pt}{3pt}{\normalfont}{}{\bfseries}{.}{.5em}
  {\llap{\textasteriskcentered\,}\thmnumber{#2}\thmnote{ #3}}
\theoremstyle{starredexercise}
\newtheorem{starredexercise}[exercise]{Exercício}
\theoremstyle{definition}

% Ambientes não numerados usados no texto
\newtheorem*{property}{Propriedade}
\newtheorem*{solution}{Solução}

% Se um enumerate for o primeiro conteúdo de um teorema/exercício (cabeçalho
% ainda pendente, \@inlabel true), o \leavevmode dispara o \everypar do amsthm
% que coloca o cabeçalho na própria linha, antes da lista iniciar.
\makeatletter
\setlist[enumerate]{before={\if@inlabel\leavevmode\fi}}
\makeatother

% Nomes em português para os ambientes
\AtBeginDocument{\renewcommand{\proofname}{\normalfont\textbf{Prova}}}
\renewcommand{\qedsymbol}{$\blacksquare$}

% Cleveref (deve vir depois de hyperref e amsthm)
\usepackage[brazilian,nameinlink]{cleveref}

\crefname{theorem}{teorema}{teoremas}
\Crefname{theorem}{Teorema}{Teoremas}
\crefname{definition}{definição}{definições}
\Crefname{definition}{Definição}{Definições}
\crefname{lemma}{lema}{lemas}
\Crefname{lemma}{Lema}{Lemas}
\crefname{proposition}{proposição}{proposições}
\Crefname{proposition}{Proposição}{Proposições}
\crefname{corollary}{corolário}{corolários}
\Crefname{corollary}{Corolário}{Corolários}
\crefname{example}{exemplo}{exemplos}
\Crefname{example}{Exemplo}{Exemplos}
\crefname{exercise}{exercício}{exercícios}
\Crefname{exercise}{Exercício}{Exercícios}
\crefname{remark}{observação}{observações}
\Crefname{remark}{Observação}{Observações}
\crefname{section}{seção}{seções}
\Crefname{section}{Seção}{Seções}

%----------------------------------------------------------------------
% MACROS ADICIONAIS
%----------------------------------------------------------------------
\let\oldemptyset\emptyset
\let\emptyset\varnothing
\graphicspath{{./}{../arquivos-aux/}}

% Seno na convenção em português
\renewcommand{\sin}{\operatorname{sen}}

\begin{document}
% --- CABEÇALHO COM O LOGO DA UnB ---
\vspace*{-0.8cm}
\begin{figure}[H]
\includegraphics{logo-mat-unb.pdf}
\end{figure}


% --- TÍTULO, AUTOR E DATA ---
\vspace*{0.3cm}
\begin{center}
{\LARGE\bf Desigualdades de Markov e Chebyshev\\[2pt] e o lema de Borel-Cantelli}
\vspace{1cm}

{\large \bf L. Cioletti}
\vspace{0.3cm}

{\footnotesize Outubro de 2026}
\end{center}

% --- RESUMO (ABSTRACT) ---
\begin{abstract}
Nestas notas, estudamos três famílias de desigualdades que relacionam
probabilidades, esperanças e momentos, e o lema de Borel-Cantelli, que decide
quando uma sequência de eventos ocorre infinitas vezes. Começamos com as
desigualdades de Markov e Chebyshev, que transformam informação sobre a média e
a variância em cotas para a probabilidade de desvios. A seguir, estudamos
convexidade: o critério da corda, a monotonicidade das inclinações secantes, a
reta de suporte e a desigualdade de Jensen, com as caracterizações
equivalentes no caso diferenciável. Com essas ferramentas, demonstramos a
desigualdade de Lyapunov, que compara momentos de ordens distintas, e
descrevemos o comportamento da curva $r\longmapsto(\mathbb{E}|X|^r)^{1/r}$. Por
fim, definimos $\limsup$ e $\liminf$ de eventos, recordamos a continuidade das
medidas de probabilidade e demonstramos as duas partes do lema de
Borel-Cantelli, com a lei 0-1 de Borel como corolário. Três aplicações fecham
os blocos: o limiar $\sqrt N$ do problema dos aniversários, o jogo honesto cujo
capital típico decai e a maior sequência de caras em $n$
lançamentos. A exposição das desigualdades de Markov e Chebyshev e do material
de convexidade segue de perto a Seção 7.5 do Capítulo 7, os problemas desse
capítulo e a Seção 8.2 do Capítulo 8 da referência \cite{GrimmettWelsh2014}. O
lema de Borel-Cantelli não consta dessa referência e segue \cite{MR1155402} e
\cite{RollaLima2026}.
\end{abstract}

% --- LICENÇA NO RODAPÉ ---
\blfootnote{%
\doclicenseImage[imagewidth=2.6em]\ 
\begingroup
\hypersetup{urlcolor=darkgray}
\href{\doclicenseURL}{%
\doclicenseName - Attribution-NonCommercial-ShareAlike 4.0 International%
}%
\endgroup
}

%----------------------------------------------------------------------
% CONTEÚDO
%----------------------------------------------------------------------
\section{Introdução}\label{sec:introducao}

Nas notas anteriores, estudamos a esperança, a variância e os momentos de
uma variável aleatória, além de sua função geradora de momentos. Para calcular
essas grandezas, partimos da distribuição da variável. Suponha agora que
conhecemos apenas sua média e sua variância. O que podemos dizer sobre a
probabilidade de a variável assumir valores em um dado intervalo? Em geral,
essas informações não permitem calcular a probabilidade procurada, pois
distribuições diferentes podem ter a mesma média e a mesma variância. Ainda
assim, podemos usá-las para obter estimativas que valem para todas essas
distribuições.

Para obter uma primeira estimativa, vamos considerar uma variável aleatória
$X$ não negativa e com esperança finita. Se $X$ assume valores maiores ou
iguais a $t>0$ com probabilidade $p$, esses valores contribuem com pelo menos
$tp$ para a esperança. Como os demais valores são não negativos, devemos ter
$tp\leqslant\mathbb{E}[X]$. Portanto,
$\mathbb{P}(X\geqslant t)\leqslant\mathbb{E}[X]/t$. Esta é a desigualdade
de Markov, que estudaremos na \Cref{sec:markov-chebyshev}. Note que o argumento
usa apenas a não negatividade de $X$ e o valor de sua esperança. Podemos,
assim, aplicá-lo a outras variáveis não negativas, em particular ao quadrado
do desvio de uma variável em relação à sua média.

O quadrado desse desvio tem esperança igual à variância. Aplicando Markov a
$(X-\mathbb{E}[X])^2$, obteremos a desigualdade de Chebyshev, que permite
estimar a probabilidade de $X$ se afastar de sua média quando a variância é
finita. Veremos uma aplicação dessas duas desigualdades no problema dos
aniversários. Se há $N$ dias possíveis para os aniversários, supostos
independentes e uniformemente distribuídos, queremos saber quantas pessoas
são necessárias para que uma coincidência se torne provável. Contando o
número de pares que fazem aniversário no mesmo dia e calculando sua média e
sua variância, mostraremos que a mudança de comportamento ocorre quando o
número de pessoas é da ordem de $\sqrt N$. Esse exemplo mostra como a escolha
da variável aleatória pode tornar úteis informações sobre apenas dois de
seus momentos.

O uso dos momentos nessas estimativas também nos leva a perguntar como
momentos de ordens distintas se relacionam. Por exemplo, sabemos que
$(\mathbb{E}[X])^2\leqslant\mathbb{E}[X^2]$, pois a variância é não
negativa. Podemos interpretar essa desigualdade como uma comparação entre
elevar a média ao quadrado e tomar a média dos quadrados. Na
\Cref{sec:convexidade}, veremos que essa comparação se estende a uma classe
de funções que inclui a função quadrática: as funções convexas. Para
entender essa extensão, começaremos pela geometria de seus gráficos.
Estudaremos as cordas e as inclinações das secantes até obter uma reta de
suporte, isto é, uma reta que passa por um ponto do gráfico e fica abaixo
dele. Essa reta nos permitirá demonstrar a desigualdade de Jensen,
$g(\mathbb{E}[X])\leqslant\mathbb{E}[g(X)]$, sob as hipóteses apropriadas.

Aplicando Jensen a funções potência, obteremos na \Cref{sec:lyapunov} a
desigualdade de Lyapunov, que compara as quantidades
$(\mathbb{E}[|X|^r])^{1/r}$ para diferentes valores de $r>0$. Em particular,
veremos que a finitude de um momento absoluto implica a finitude dos momentos
absolutos de ordens menores. A comparação entre a função da média e a média
da função também será útil para estudar um jogo em que o capital é
multiplicado, a cada rodada, por um fator aleatório de média igual a $1$.
Embora a esperança do capital permaneça constante, mostraremos que a
probabilidade de ele ficar abaixo de qualquer valor positivo fixado tende a
$1$ com o número de rodadas. O logaritmo transforma os produtos em somas, e a
desigualdade de Jensen ajuda a entender por que a esperança do logaritmo do
capital decresce mesmo quando a esperança do capital não muda.

Ao acompanhar o capital ao longo das rodadas, passamos a considerar uma
sequência de variáveis aleatórias. Nesse tipo de situação, além de estimar
uma probabilidade em cada etapa, podemos querer saber o que acontece ao
longo de toda a sequência. Por exemplo, se $A_n$ é o evento de ocorrer um
desvio na etapa $n$, será que esses desvios deixam de ocorrer a partir de
algum momento ou continuam ocorrendo infinitas vezes? Mesmo que
$\mathbb{P}(A_n)$ tenda a zero, essa informação, sozinha, não decide a
questão. Precisamos levar em conta o conjunto das probabilidades e a
dependência entre os eventos.

É dessa questão que trata o lema de Borel-Cantelli, estudado na
\Cref{sec:borel-cantelli}. Se a série $\sum_n\mathbb{P}(A_n)$ converge,
então, com probabilidade $1$, somente um número finito dos eventos $A_n$
ocorre. Se os eventos são independentes e a série diverge, então, com
probabilidade $1$, infinitos deles ocorrem. Para formular e demonstrar essas
afirmações, introduziremos os limites superior e inferior de eventos e
recordaremos a continuidade das medidas de probabilidade. Ao final,
aplicaremos essas ferramentas à maior sequência de caras em $n$ lançamentos
independentes de uma moeda honesta, mostrando que seu comprimento, dividido
por $\log_2 n$, tende a $1$ quase certamente. A demonstração partirá de
estimativas para as probabilidades de sequências de caras muito longas ou
muito curtas. Comecemos, portanto, pelas ferramentas que permitem obter
estimativas de probabilidades a partir de esperanças.

%----------------------------------------------------------------------
\section{As Desigualdades de Markov e Chebyshev}\label{sec:markov-chebyshev}

\subsection{A Desigualdade de Markov}\label{subsec:markov}

A pergunta que orienta esta seção é a seguinte: se sabemos que
$\mathbb{E}[X]$ é finito e que $X$ não assume valores negativos, o que podemos
dizer sobre a probabilidade de $X$ ser grande? A cauda direita de $X$ é a
função $t\longmapsto\mathbb{P}(X\geqslant t)$, definida para $t>0$. A
desigualdade de Markov diz que essa função não pode decair mais devagar do que
$1/t$.

\begin{theorem}[Desigualdade de Markov]\label{thm:markov}
Seja $X$ uma variável aleatória não negativa. Então, para cada $t>0$,
\begin{equation}\label{eq:markov}
\mathbb{P}(X\geqslant t)\leqslant\frac{\mathbb{E}[X]}{t}.
\end{equation}
\end{theorem}

\begin{proof}
Fixe $t>0$. Os eventos $\{X\geqslant t\}$ e $\{X<t\}$ são disjuntos e a união
deles é $\Omega$, de modo que formam uma partição de $\Omega$. A esperança é
aditiva nessa decomposição, e a forma de escrever isso é multiplicar $X$ pela
indicadora de cada uma das partes:
\begin{align*}
\mathbb{E}[X]
&=
\mathbb{E}\Bigl[X\Bigl(\mathds{1}_{\{X\geqslant t\}}
+\mathds{1}_{\{X<t\}}\Bigr)\Bigr]
\\[0.3cm]
&=
\mathbb{E}\bigl[X\,\mathds{1}_{\{X\geqslant t\}}\bigr]
+
\mathbb{E}\bigl[X\,\mathds{1}_{\{X<t\}}\bigr].
\end{align*}
A segunda parcela é não negativa, porque $X\geqslant0$. Descartá-la só pode
diminuir o lado direito, logo
\[
\mathbb{E}[X]\geqslant\mathbb{E}\bigl[X\,\mathds{1}_{\{X\geqslant t\}}\bigr].
\]
No evento $\{X\geqslant t\}$, a variável $X$ é limitada inferiormente por $t$.
Substituir $X$ por $t$ nessa parcela só pode diminuir a esperança, donde
\[
\mathbb{E}\bigl[X\,\mathds{1}_{\{X\geqslant t\}}\bigr]
\geqslant
\mathbb{E}\bigl[t\,\mathds{1}_{\{X\geqslant t\}}\bigr]
=
t\,\mathbb{P}(X\geqslant t).
\]
Reunindo as duas desigualdades, obtemos
$\mathbb{E}[X]\geqslant t\,\mathbb{P}(X\geqslant t)$. Como $t>0$, podemos
dividir ambos lados da desigualdade por $t$ e concluir a validade de \eqref{eq:markov}.
\end{proof}

A cota é ótima no sentido de que não se pode melhorar a constante. Se $X$
assume o valor $t$ com probabilidade $p$ e o valor $0$ com probabilidade
$1-p$, então $\mathbb{E}[X]=pt$ e $\mathbb{P}(X\geqslant t)=p$, de modo que
\eqref{eq:markov} vale com igualdade. O que a desigualdade não controla é a
forma da cauda: duas variáveis com a mesma média podem ter caudas muito
diferentes, e Markov dá a mesma cota para ambas.

A prova não usou praticamente nenhuma propriedade específica da distribuição da variável
aleatória $X$, apenas a positividade e a
linearidade da esperança. Por isso ela se estende sem mudança a qualquer
potência ou função crescente de $X$. Se $g\colon[0,\infty)\to[0,\infty)$ é não
decrescente e $g(t)>0$, então o evento $\{X\geqslant t\}$ está contido em
$\{g(X)\geqslant g(t)\}$, e a desigualdade de Markov aplicada a $g(X)$ dá
\begin{equation}\label{eq:markov-generalizada}
\mathbb{P}(X\geqslant t)
\leqslant
\mathbb{P}\bigl(g(X)\geqslant g(t)\bigr)
\leqslant
\frac{\mathbb{E}[g(X)]}{g(t)}.
\end{equation}
A desigualdade de Chebyshev é o caso $g(x)=x^2$, e é essa escolha que
introduz a variância.

\subsection{A Desigualdade de Chebyshev}\label{subsec:chebyshev}

Se $X$ tem média $\mu$ e variância finita, então a variável aleatória $(X-\mu)^2$ é não
negativa e sua esperança é $\operatorname{Var}(X)$. Aplicando a desigualdade de
Markov a essa variável, obtemos uma cota para a probabilidade de $X$ se
afastar de $\mu$.

\begin{theorem}[Desigualdade de Chebyshev]\label{thm:chebyshev}
Seja $X$ uma variável aleatória com $\mathbb{E}[X^2]<\infty$. Então, para cada
$\varepsilon>0$,
\begin{equation}\label{eq:chebyshev}
\mathbb{P}\bigl(|X-\mathbb{E}[X]|\geqslant\varepsilon\bigr)
\leqslant
\frac{\operatorname{Var}(X)}{\varepsilon^2}.
\end{equation}
\end{theorem}

\begin{proof}
Seja $\mu=\mathbb{E}[X]$. A função $u\longmapsto u^2$ é crescente em $[0,\infty)$
e ambos os membros na desigualdade $|X-\mu|\geqslant\varepsilon$ são não negativos, logo
\[
\{|X-\mu|\geqslant\varepsilon\}
=
\{(X-\mu)^2\geqslant\varepsilon^2\}.
\]
Aplicando o \Cref{thm:markov} à variável não negativa $(X-\mu)^2$, com $t=\varepsilon^2$, obtemos
\[
\mathbb{P}\bigl((X-\mu)^2\geqslant\varepsilon^2\bigr)
\leqslant
\frac{\mathbb{E}\bigl[(X-\mu)^2\bigr]}{\varepsilon^2}
=
\frac{\operatorname{Var}(X)}{\varepsilon^2}.
\]
Isso prova \eqref{eq:chebyshev}.
\end{proof}

Equivalentemente, tomando o complementar,
\begin{equation}\label{eq:chebyshev-complementar}
\mathbb{P}\bigl(|X-\mathbb{E}[X]|<\varepsilon\bigr)
\geqslant
1-\frac{\operatorname{Var}(X)}{\varepsilon^2}.
\end{equation}
A forma \eqref{eq:chebyshev-complementar} é a que costuma ser usada: ela diz
que, com probabilidade alta, $X$ fica em um intervalo centrado na média, e que
o comprimento desse intervalo precisa crescer como o desvio padrão para que a
probabilidade se aproxime de $1$.

A mesma prova, com $|X-\mu|^p$ no lugar de $(X-\mu)^2$, dá
\[
\mathbb{P}\bigl(|X-\mu|\geqslant\varepsilon\bigr)
\leqslant
\frac{\mathbb{E}\bigl[|X-\mu|^p\bigr]}{\varepsilon^p},
\]
para todo $p>0$ tal que o momento exista. Quanto maior o momento disponível,
mais rápido é o decaimento na cota superior. A \Cref{sec:lyapunov} mostra que os momentos
de ordens maiores são os mais exigentes, e que a existência deles é uma
condição mais forte.

\begin{remark}\label{rem:chebyshev-nao-e-otima}
A desigualdade de Chebyshev é geral, e é justamente por isso que ela não é
ótima em exemplos concretos. Para $X\sim\mathrm{Bin}(n,p)$, por exemplo, a cota
para $\mathbb{P}(|X-np|\geqslant\varepsilon)$ é da ordem de $1/\varepsilon^2$,
enquanto a distribuição real decai de forma gaussiana. O valor da desigualdade
está em não exigir nada sobre a forma da distribuição.
\end{remark}

\subsection{A Mediana e uma Cota para a Média}\label{subsec:mediana}

A desigualdade de Markov também compara a média com a mediana, e essa
comparação é útil porque a mediana sempre existe, mesmo quando a média ou valor
esperado não está bem-definido.

\begin{definition}[Mediana]\label{def:mediana}
Seja $X$ uma variável aleatória. Um número real $m$ é uma \emph{mediana} de
$X$ se
\[
\mathbb{P}(X<m)\leqslant\frac{1}{2}\leqslant\mathbb{P}(X\leqslant m).
\]
\end{definition}

Toda variável aleatória possui pelo menos uma mediana. De fato, seja $F$ a
função de distribuição de $X$ e defina
\[
m=\inf\{x\in\mathbb{R}:F(x)\geqslant1/2\}.
\]
% Funções de distribuição discretas: mediana única e intervalo de medianas.
\begin{figure}[H]
\centering
\begin{minipage}[t]{0.48\textwidth}
\centering
\begin{tikzpicture}[x=1.35cm, y=3.6cm, font=\scriptsize]
  % Eixos e níveis de referência.
  \draw[-stealth] (-0.25,0) -- (4.35,0) node[below right] {$x$};
  \draw[-stealth] (0,-0.08) -- (0,1.18) node[above left] {$F(x)$};
  \draw[gray,dashed] (0,0.5) -- (4.1,0.5);
  \draw[gray,dashed] (0,1) -- (4.1,1);
  \node[left] at (0,0.5) {$\frac12$};
  \node[left] at (0,1) {$1$};
  \node[below left] at (0,0) {$0$};

  % Massas 1/4, 1/2 e 1/4, nas abscissas 0.7, 2 e 3.3.
  \draw[thick] (-0.25,0) -- (0.7,0);
  \draw[thick] (0.7,0.25) -- (2,0.25);
  \draw[thick] (2,0.75) -- (3.3,0.75);
  \draw[thick] (3.3,1) -- (4.1,1);
  \draw[gray,dashed] (0.7,0) -- (0.7,0.25);
  \draw[teal,dashed] (2,0) -- (2,0.75);
  \draw[gray,dashed] (3.3,0) -- (3.3,1);
  \foreach \p in {(0.7,0.25),(3.3,1)}
    \fill \p circle (1.6pt);
  \foreach \p in {(0.7,0),(2,0.25),(3.3,0.75)}
    \draw[teal,fill=white,thick] \p circle (1.7pt);
  \fill[teal] (2,0.75) circle (1.7pt);
  \node[below left] at (2,0.25) {$F(m^-)\!<\!\frac12$};
  \node[above right] at (2,0.75) {$F(m)\!>\!\frac12$};

  % A única mediana, destacada no eixo.
  \fill[teal] (2,0) circle (2pt);
  \node[below,teal] at (2,-0.015) {$m$};
  \node at (2,-0.23) {(a) Mediana única};
\end{tikzpicture}
\end{minipage}\hfill
\begin{minipage}[t]{0.48\textwidth}
\centering
\begin{tikzpicture}[x=1.35cm, y=3.6cm, font=\scriptsize]
  \draw[-stealth] (-0.25,0) -- (4.35,0) node[below right] {$x$};
  \draw[-stealth] (0,-0.08) -- (0,1.18) node[above left] {$F(x)$};
  \draw[gray,dashed] (0,0.5) -- (4.1,0.5);
  \draw[gray,dashed] (0,1) -- (4.1,1);
  \node[left] at (0,0.5) {$\frac12$};
  \node[left] at (0,1) {$1$};
  \node[below left] at (0,0) {$0$};

  % Probabilidade 1/2 em a e 1/2 em b.
  \draw[thick] (-0.25,0) -- (1.2,0);
  \draw[teal,thick] (1.2,0.5) -- (2.9,0.5);
  \draw[thick] (2.9,1) -- (4.1,1);
  \draw[gray,dashed] (1.2,0) -- (1.2,0.5);
  \draw[gray,dashed] (2.9,0) -- (2.9,1);
  \foreach \p in {(1.2,0),(2.9,0.5)}
    \draw[teal,fill=white,thick] \p circle (1.7pt);
  \fill[teal] (1.2,0.5) circle (1.7pt);
  \fill (2.9,1) circle (1.6pt);

  % O extremo b também é mediana, embora F(b)=1.
  \draw[teal,very thick] (1.2,0) -- (2.9,0);
  \fill[teal] (1.2,0) circle (2pt);
  \fill[teal] (2.9,0) circle (2pt);
  \node[below,teal] at (1.2,-0.015) {$a=m$};
  \node[below,teal] at (2.9,-0.015) {$b$};
  \node[above,teal] at (2.05,0.02) {medianas};
  \node at (2,-0.23) {(b) Intervalo de medianas};
\end{tikzpicture}
\end{minipage}
\caption{A mediana separa a distribuição de modo que a probabilidade de
ficar estritamente de cada lado dela é no máximo $1/2$. À esquerda, somente
o número $m$ satisfaz essa condição: $F(m^-)=\mathbb{P}(X<m)<1/2<F(m)$.
À direita, $X$ assume os valores $a$ e $b$ com probabilidade $1/2$ cada,
e todo ponto de $[a,b]$ é mediana. O extremo $b$ também está incluído,
pois $\mathbb{P}(X<b)=1/2$, embora $F(b)=1$. A definição por ínfimo
seleciona a menor mediana, que nesse caso é $m=a$.}
\label{fig:mediana-unicidade}
\end{figure}

Vamos justificar com cuidado por que esse ínfimo existe e por que ele
fornece uma mediana. Para simplificar a notação, escrevemos
\[
S=\{x\in\mathbb{R}:F(x)\geqslant1/2\}.
\]
O conjunto $S$ reúne os pontos nos quais a probabilidade acumulada já
alcançou pelo menos $1/2$. Como $F(x)$ tende a $1$ quando $x$ tende a
$+\infty$, podemos escolher um número $b$ suficientemente grande para que
$F(b)>1/2$. Assim, $b\in S$ e $S$ não é vazio. Na outra direção,
$F(x)$ tende a $0$ quando $x$ tende a $-\infty$. Podemos, portanto,
escolher um número $a$ tal que $F(a)<1/2$. Pela monotonicidade de $F$,
\[
x\leqslant a\quad\Longrightarrow\quad F(x)\leqslant F(a)<1/2.
\]
Nenhum ponto à esquerda de $a$ pertence a $S$. Logo $a$ é uma cota inferior
de $S$, isto é, um número menor ou igual a todos os elementos de $S$.
Como $S$ é não vazio e possui uma cota inferior, a propriedade de
completude dos números reais garante que $m=\inf S$ é um número real.

O ínfimo é a \emph{maior das cotas inferiores} de $S$. Vamos traduzir essa
descrição em duas informações que usaremos na prova. Primeiro, como $m$ é
uma cota inferior, todo $y\in S$ satisfaz $m\leqslant y$. Em particular,
se $x<m$, então $x\notin S$, ou seja,
\[
x<m\quad\Longrightarrow\quad F(x)<1/2.
\]
Segundo, como $m$ é a maior cota inferior, nenhum número maior que $m$
pode ser uma cota inferior de $S$. Assim, dado qualquer $\varepsilon>0$,
existe algum $y\in S$ tal que
\[
m\leqslant y<m+\varepsilon.
\]
De fato, se não existisse tal $y$, todos os elementos de $S$ seriam maiores
ou iguais a $m+\varepsilon$. Esse número seria uma cota inferior de $S$
maior que $m$, contradizendo a escolha de $m$.

Essa segunda informação permite controlar $F$ à direita de $m$. Como
$y\in S$ e $y<m+\varepsilon$, a monotonicidade dá
\[
F(m+\varepsilon)\geqslant F(y)\geqslant1/2.
\]
Portanto, à esquerda de $m$ temos $F(x)<1/2$ e à direita temos
$F(x)\geqslant1/2$. Falta verificar o que acontece no próprio ponto $m$.
Essa verificação é necessária: o ínfimo de um conjunto nem sempre pertence
ao conjunto. Por exemplo, o intervalo $(0,1)$ tem ínfimo $0$, mas não
contém $0$.

% O papel das duas propriedades do ínfimo antes de usar continuidade.
\begin{figure}[H]
\centering
\begin{tikzpicture}[x=1.45cm,y=1cm,font=\small]
  \draw[-stealth] (-4.4,0) -- (4.4,0) node[below right] {$x$};
  % Destacamos apenas os pontos estritamente à direita de m:
  % a pertinência de m a S ainda será demonstrada no texto.
  \draw[teal,very thick] (0.04,0) -- (4.1,0);
  \draw (0,-0.1) -- (0,0.1);
  \node[below] at (0,-0.1) {$m=\inf S$};
  \draw (2.6,-0.08) -- (2.6,0.08);
  \node[below] at (2.6,-0.1) {$m+\varepsilon$};
  \fill[teal] (1.1,0) circle (2pt);
  \node[below,teal] at (1.1,-0.1) {$y\in S$};
  \node at (-2.2,0.7) {$x<m\;\Longrightarrow\;F(x)<\frac12$};
  \node[teal] at (2.2,0.7)
    {$x>m\;\Longrightarrow\;F(x)\geqslant\frac12$};
  \draw[gray,dashed] (0,0.18) -- (0,1.45);
  \draw[-stealth,teal] (2.6,1.25) -- (0.16,1.25);
  \node[above,teal,font=\scriptsize] at (1.4,1.3)
    {aproximação pela direita};
  \draw[-stealth] (-2.6,1.25) -- (-0.16,1.25);
  \node[above,font=\scriptsize] at (-1.4,1.3)
    {aproximação pela esquerda};
\end{tikzpicture}
\caption{As duas propriedades do ínfimo localizam os pontos em relação a
$m$. Nenhum ponto menor que $m$ pertence a $S$; para cada
$\varepsilon>0$, há um ponto $y\in S$ entre $m$ e $m+\varepsilon$
(possivelmente $y=m$). A monotonicidade de $F$ garante então que todos os
pontos maiores que $m$ pertencem a $S$. A inclusão do próprio $m$ será
justificada usando a continuidade à direita.}
\label{fig:infimo-mediana}
\end{figure}


Para decidir o que acontece em $m$, vamos aproximá-lo por pontos à sua
direita, onde já sabemos que $F$ é pelo menos $1/2$. Tomando
$\varepsilon=1/n$, obtemos
\[
F(m+1/n)\geqslant1/2\qquad\text{para todo }n\geqslant1.
\]
Quando $n$ cresce, $m+1/n$ se aproxima de $m$ pela direita. É justamente
para essa aproximação que podemos usar a continuidade à direita da função
de distribuição:
\[
\mathbb{P}(X\leqslant m)=F(m)
=\lim_{n\to\infty}F(m+1/n)\geqslant1/2.
\]
Obtivemos, assim, uma das desigualdades exigidas na definição de mediana.

Para obter a outra, precisamos calcular $\mathbb{P}(X<m)$. Agora a
aproximação deve ser feita pela esquerda, pois queremos excluir o valor
$m$. Cada ponto $m-1/n$ é menor que $m$ e, portanto,
\[
\mathbb{P}(X\leqslant m-1/n)=F(m-1/n)<1/2.
\]
Os eventos $\{X\leqslant m-1/n\}$ crescem com $n$ e sua união é
$\{X<m\}$. Para entender a igualdade da união, note que um valor
estritamente menor que $m$ tem uma distância positiva até $m$.
Quando $1/n$ fica menor que essa distância, esse valor já pertence ao
evento $\{X\leqslant m-1/n\}$. Por outro lado, o valor $m$ não pertence
a nenhum desses eventos. Pela continuidade da probabilidade para
sequências crescentes, que recordaremos na \Cref{subsec:continuidade},
\[
\mathbb{P}(X<m)
=\lim_{n\to\infty}\mathbb{P}(X\leqslant m-1/n)
=\lim_{n\to\infty}F(m-1/n)\leqslant1/2.
\]
Note que a desigualdade no limite pode deixar de ser estrita: números
menores que $1/2$ podem se aproximar de $1/2$. Reunindo as duas conclusões,
\[
\mathbb{P}(X<m)\leqslant\frac12
\leqslant\mathbb{P}(X\leqslant m).
\]
Assim, o número $m$ escolhido é uma mediana de $X$. A aproximação pela
direita garantiu que a probabilidade acumulada em $m$ já alcançou $1/2$;
a aproximação pela esquerda garantiu que, antes de $m$, ela não ultrapassa
$1/2$.

\begin{proposition}\label{prop:mediana-markov}
Seja $X$ uma variável aleatória não negativa com $\mathbb{E}[X]<\infty$. Se
$m$ é uma mediana de $X$, então $m\leqslant2\mathbb{E}[X]$.
\end{proposition}

\begin{proof}
Como $m$ é mediana, $\mathbb{P}(X\geqslant m)\geqslant1/2$. Pelo
\Cref{thm:markov},
\[
\frac{1}{2}\leqslant\mathbb{P}(X\geqslant m)\leqslant\frac{\mathbb{E}[X]}{m}.
\]
Se $m>0$, basta multiplicar a última desigualdade por $2m$.
Se $m\leqslant0$, a conclusão é imediata, porque $\mathbb{E}[X]\geqslant0$.
\end{proof}

A proposição diz que a mediana de uma variável não negativa não pode estar
muito à direita da média. A recíproca não vale: a média pode estar muito à
direita da mediana, e é isso que acontece com variáveis de cauda pesada. O
\Cref{exercise:mediana-igualdade} pede o exemplo em que a cota é atingida.

\subsection{Aplicação: O Limiar do Problema dos Aniversários}\label{subsec:aniversarios}

Em uma sala com $n$ pessoas, qual é a probabilidade de que duas façam
aniversário no mesmo dia? A resposta exata depende de uma contagem, mas a
ordem de grandeza do número de pessoas necessário para que a coincidência se
torne provável pode ser obtida com as desigualdades de Markov e Chebyshev,
sem calcular explicitamente a distribuição. Nosso objetivo será determinar
essa escala: queremos entender como o número de pessoas precisa crescer
quando o número $N$ de dias possíveis aumenta. Esse é um exemplo de uma
pergunta assintótica, na qual procuramos identificar o comportamento do
problema para valores grandes dos parâmetros.

Suponha que os aniversários sejam independentes e uniformemente distribuídos
em $N$ dias, e que $N$ seja grande. Vamos supor $n\geqslant2$.
Para detectar uma coincidência, basta que algum par de pessoas faça
aniversário no mesmo dia. Isso sugere contar os pares coincidentes.
Numeramos as pessoas de $1$ a $n$ e, para
cada par $\{i,j\}$ com $i<j$, definimos
\[
I_{ij}=\mathds{1}_{\{i\text{ e }j\text{ fazem aniversário no mesmo dia}\}},
\qquad
X=\sum_{1\leqslant i<j\leqslant n}I_{ij}.
\]
Assim $X$ é a variável aleatória que conta o número de pares de pessoas, na sala,
que fazem aniversário no mesmo dia. Não estamos contando pessoas: se três
delas fazem aniversário no mesmo dia, por exemplo, elas contribuem com três
pares para $X$. O que importa é que existe tal par se, e somente se, $X\geqslant1$. Logo o evento de interesse é
$\{X\geqslant1\}$. Para cada par de pessoas, qualquer que seja o dia do
aniversário da primeira, a segunda tem probabilidade $1/N$ de fazer
aniversário nesse mesmo dia, pela independência e pela distribuição uniforme
dos aniversários. Assim, $\mathbb{E}[I_{ij}]=1/N$ e, pela linearidade da
esperança, somando sobre os $\binom{n}{2}$ pares, obtemos
\[
\mathbb{E}[X]=\binom{n}{2}\frac{1}{N}.
\]
Pela desigualdade de Markov,
\begin{equation}\label{eq:aniversarios-markov}
\mathbb{P}(X\geqslant1)
\leqslant \frac{\mathbb{E}[X]}{1}
=
\binom{n}{2}\frac{1}{N}
=
\frac{n(n-1)}{2} \frac{1}{N}
\approx\frac{n^2}{2N}.
\end{equation}
O termo $n^2/N$ já sugere a escala $\sqrt N$: ele é pequeno quando $n$ é
pequeno em relação a $\sqrt N$. Mais precisamente, se fazemos $N$ crescer
e escolhemos $n=n(N)$ de modo que $n/\sqrt N\to0$, então
\[
\mathbb{P}(X\geqslant1)
\leqslant\frac{n(n-1)}{2N}
\leqslant\frac12\left(\frac{n}{\sqrt N}\right)^2
\longrightarrow0.
\]
É esse o sentido da notação $n\ll\sqrt N$. Para concluir que $\sqrt N$
é a escala da mudança de comportamento, falta estudar o outro lado:
o que acontece quando $n$ é grande em relação a $\sqrt N$?

Nesse caso, esperamos muitas coincidências, pois $\mathbb{E}[X]$ é grande.
Uma esperança grande, sozinha, não garante que $X$ seja positivo com alta
probabilidade: ela poderia ser sustentada por valores muito grandes e muito
raros de $X$. Para afastar essa possibilidade, vamos calcular a variância e
usar Chebyshev para controlar os desvios de $X$ em relação à sua média.

Comecemos pela variância de cada indicadora. Como $I_{ij}$ assume apenas
os valores $0$ e $1$, temos $I_{ij}^2=I_{ij}$ e, portanto,
\[
\operatorname{Var}(I_{ij})
=\mathbb{E}[I_{ij}^2]-(\mathbb{E}[I_{ij}])^2
=\frac1N-\frac1{N^2}
=\frac1N\left(1-\frac1N\right).
\]
Para obter a variância da soma, precisamos também das covariâncias entre
indicadoras distintas. Não podemos simplesmente supor que elas sejam
independentes: dois pares podem compartilhar uma pessoa. Vamos verificar
os dois casos, representados na \Cref{fig:aniversarios-pares}.

\begin{figure}[H]
\centering
\begin{tikzpicture}[
  x=1cm,y=1cm,font=\small,
  pessoa/.style={circle,draw,fill=white,minimum size=7mm,inner sep=0pt},
  coincidencia/.style={thick,teal}]
  % Arestas representam as coincidências exigidas por cada indicadora.
  \node[pessoa] (a1) at (0,0) {$1$};
  \node[pessoa] (a2) at (2,0) {$2$};
  \node[pessoa] (a3) at (0,-1.3) {$3$};
  \node[pessoa] (a4) at (2,-1.3) {$4$};
  \draw[coincidencia] (a1) -- node[above] {$I_{12}=1$} (a2);
  \draw[coincidencia] (a3) -- node[above] {$I_{34}=1$} (a4);
  \node at (1,-2.35) {Pares sem pessoas em comum};

  \node[pessoa] (b1) at (7,0.3) {$1$};
  \node[pessoa] (b2) at (5.7,-1.3) {$2$};
  \node[pessoa] (b3) at (8.3,-1.3) {$3$};
  \draw[coincidencia] (b1) -- node[above left] {$I_{12}=1$} (b2);
  \draw[coincidencia] (b1) -- node[above right] {$I_{13}=1$} (b3);
  \draw[gray,dashed] (b2) -- node[below=7pt,font=\scriptsize,text=black!65]
    {também $I_{23}=1$} (b3);
  \node at (7,-2.35) {Pares com uma pessoa em comum};
\end{tikzpicture}
\caption{Cada círculo representa uma pessoa e cada segmento contínuo exige
que seus extremos façam aniversário no mesmo dia. À esquerda, as duas
coincidências envolvem quatro pessoas. À direita, exigir as duas
coincidências faz com que as três pessoas tenham o mesmo aniversário;
a terceira coincidência, tracejada, é uma consequência das outras duas.}
\label{fig:aniversarios-pares}
\end{figure}


Se os pares não têm pessoas em comum, como $\{1,2\}$ e $\{3,4\}$,
as duas coincidências dependem de grupos distintos de aniversários
independentes. Logo
\[
\mathbb{P}(I_{12}=1,\,I_{34}=1)
=\mathbb{P}(I_{12}=1)\mathbb{P}(I_{34}=1)
=\frac1{N^2}.
\]
Se os pares compartilham uma pessoa, como $\{1,2\}$ e $\{1,3\}$,
as duas coincidências ocorrem exatamente quando as três pessoas fazem
aniversário no mesmo dia. Para cada dia fixado, essa probabilidade é
$(1/N)^3$, pela independência dos três aniversários. Como há $N$ dias
possíveis e esses casos são disjuntos,
\[
\mathbb{P}(I_{12}=1,\,I_{13}=1)
=N\left(\frac1N\right)^3
=\frac1{N^2}
=\mathbb{P}(I_{12}=1)\mathbb{P}(I_{13}=1).
\]
Assim, nos dois casos a probabilidade conjunta é o produto das
probabilidades. Como as variáveis são indicadoras, isso mostra que são
duas a duas independentes. Em particular, para quaisquer dois pares
distintos, o produto das indicadoras tem esperança $1/N^2$ e sua
covariância é $1/N^2-(1/N)(1/N)=0$. A fórmula da variância de uma soma
fica então reduzida à soma das variâncias:
\[
\operatorname{Var}(X)
=
\sum_{i<j}\operatorname{Var}(I_{ij})
=
\binom{n}{2}\frac{1}{N}\Bigl(1-\frac{1}{N}\Bigr).
\]
Vale observar que a independência é apenas duas a duas. Por exemplo,
$I_{12}=I_{13}=1$ força $I_{23}=1$. Isso não atrapalha o cálculo da
variância, no qual aparecem somente covariâncias entre duas variáveis.

Agora queremos usar essa variância para limitar a probabilidade de não
haver pares de pessoas fazendo aniversário no mesmo dia.
Por questão de simplicidade denote por $\mu\equiv \mathbb{E}[X]=\binom n2/N>0$.
Se $X=0$, então $|X-\mu|=\mu$: a ausência de coincidências corresponde
a um desvio em relação à média de tamanho igual à própria média. Portanto,
$\{X=0\}\subseteq\{|X-\mu|\geqslant\mu\}$. Essa inclusão indica a
escolha $\varepsilon=\mu$ na desigualdade de Chebyshev que vale a seguinte desigualdade
\[
\mathbb{P}(X=0)
\leqslant
\mathbb{P}\bigl(|X-\mathbb{E}[X]|\geqslant\mathbb{E}[X]\bigr)
\leqslant
\frac{\operatorname{Var}(X)}{(\mathbb{E}[X])^2}.
\]
Substituindo a média e a variância que calculamos, a razão acima fica
\[
\frac{\operatorname{Var}(X)}{(\mathbb{E}[X])^2}
=\frac{\binom n2\,\frac1N\left(1-\frac1N\right)}
       {\left(\binom n2\,\frac1N\right)^2}
=\frac{N\left(1-\frac1N\right)}{\binom n2}
=\frac{N-1}{\binom n2}.
\]
Tomando o complementar do evento $\{X=0\}$, obtemos
\begin{equation}\label{eq:aniversarios-chebyshev}
\mathbb{P}(X\geqslant1)\geqslant1-\frac{N-1}{\binom{n}{2}}
\approx1-\frac{2N}{n^2}.
\end{equation}
Se $n/\sqrt N\to\infty$, então $n\to\infty$ e
\[
\frac{N-1}{\binom n2}
=\frac{2(N-1)}{n(n-1)}
\longrightarrow0.
\]
Logo $\mathbb{P}(X\geqslant1)\to1$. Escrevemos $n\gg\sqrt N$ para
indicar esse regime.

Reunindo \eqref{eq:aniversarios-markov} e \eqref{eq:aniversarios-chebyshev},
vemos que o número de pessoas necessário para mudar o comportamento do
problema é da ordem de $\sqrt N$. Abaixo dessa escala, no sentido
$n/\sqrt N\to0$, a probabilidade de coincidência tende a $0$; acima dela,
no sentido $n/\sqrt N\to\infty$, tende a $1$. Esse é o ganho de determinar
a escala: mesmo sem uma fórmula exata, sabemos em quais regimes o evento é
raro e em quais é muito provável. Sabemos também onde procurar uma
transição entre os dois comportamentos, isto é, quando $n$ é comparável a
$\sqrt N$. As estimativas ainda não descrevem precisamente essa transição,
mas identificam a relação entre os parâmetros que governa o problema.

\bigskip


Para avaliar a precisão das cotas em um caso concreto, tome $N=365$.
Com $n=23$, Markov fornece a cota superior
$\binom{23}{2}/365=253/365\approx0{,}69$, enquanto a probabilidade exata é
aproximadamente $0{,}51$. Com $n=41$, a cota inferior de
\eqref{eq:aniversarios-chebyshev}
dá
\[
\mathbb{P}(X\geqslant1)\geqslant1-\frac{364}{820}\approx0{,}556,
\]
enquanto a probabilidade exata é aproximadamente $0{,}903$. As cotas
localizam a escala, mas podem ser pouco precisas para valores específicos
de $n$ e $N$.

Para comparar com as cotas, vamos obter a fórmula exata. Se os aniversários
das primeiras $k$ pessoas são distintos, eles ocupam $k$ dos $N$ dias.
A probabilidade de a pessoa seguinte evitar esses dias é $(N-k)/N$.
Para multiplicar essas probabilidades condicionais, usamos a
\emph{regra da multiplicação},
\[
\mathbb{P}(A\cap B)=\mathbb{P}(A)\,\mathbb{P}(B\mid A),
\qquad \mathbb{P}(A)>0.
\]
Sua aplicação sucessiva fornece a seguinte igualdade
\[
\mathbb{P}\left(\bigcap_{j=1}^n A_j\right)
=\mathbb{P}(A_1)\prod_{j=2}^n
\mathbb{P}\left(A_j\,\middle|\,\bigcap_{i=1}^{j-1}A_i\right),
\]
desde que os eventos nos quais condicionamos tenham probabilidade positiva.
Essa fórmula foi demonstrada na Proposição 7 da Subseção 2.3,
``A Regra da Multiplicação'', das notas
\href{https://mat.unb.br/~cioletti/ensino/independencia.pdf}{\emph{Probabilidade Condicional e
Independência}}.

Aqui, se $D_k$ é o evento de que os aniversários das primeiras $k$ pessoas
são distintos, temos $D_{k+1}\subseteq D_k$. Como mencionado acima
segue da regra da multiplicação que
\[
\mathbb{P}(D_{k+1})=\mathbb{P}(D_k)\mathbb{P}(D_{k+1}\mid D_k),
\]
com $\mathbb{P}(D_1)=1$ e
$\mathbb{P}(D_{k+1}\mid D_k)=(N-k)/N$. Aplicando essa identidade
sucessivamente, obtemos, para $n\leqslant N$,
\[
\mathbb{P}(X=0)
=1\cdot\frac{N-1}{N}\cdot\frac{N-2}{N}\cdots\frac{N-n+1}{N}
=\prod_{k=0}^{n-1}\Bigl(1-\frac{k}{N}\Bigr).
\]
Se $n>N$, necessariamente há uma coincidência. Para $n\leqslant N$, a
probabilidade procurada é $1$ menos o produto acima.


\bigskip

O produto também permite enxergar a escala $\sqrt N$ por outro caminho.
Quando $n$ é pequeno em relação a $N$, cada $k/N$ é pequeno. Usando
$\log(1-u)\approx-u$, obtemos
\[
\log\mathbb{P}(X=0)
=\sum_{k=0}^{n-1}\log\left(1-\frac{k}{N}\right)
\approx-\frac1N\sum_{k=0}^{n-1}k
=-\frac{n(n-1)}{2N}.
\]
Assim, na região da transição,
\[
\mathbb{P}(X\geqslant1)
\approx1-\exp\left(-\frac{n(n-1)}{2N}\right).
\]
Para essa probabilidade ficar próxima de $1/2$, devemos ter
$\exp(-n(n-1)/(2N))\approx1/2$, ou seja,
$n(n-1)\approx2N\log2$. Para $N$ grande, isso corresponde a
$n\approx\sqrt{2\log2}\,\sqrt N$. A contagem exata permite, portanto,
refinar a escala com uma constante para o nível de probabilidade $1/2$.
Markov e Chebyshev já haviam identificado a escala $\sqrt N$ usando apenas
a média e a variância do número de pares coincidentes.

\bigskip

\section{Convexidade e a Desigualdade de Jensen}\label{sec:convexidade}

Na seção anterior, usamos a média e a variância para estimar probabilidades.
Vamos agora estudar uma relação entre os próprios momentos. Já sabemos que
$(\mathbb{E}[X])^2\leqslant\mathbb{E}[X^2]$, pois a variância é não
negativa. Nessa desigualdade, comparamos duas operações: primeiro tomar a
média e depois elevar ao quadrado, ou primeiro elevar ao quadrado e depois
tomar a média. Que outras funções, no lugar do quadrado, permitem uma
comparação semelhante?

Para responder, vamos olhar para o gráfico da função quadrática. Entre dois
pontos quaisquer de seu gráfico, a corda que os une fica acima do gráfico.
Essa propriedade define as funções convexas, como ilustra a
\Cref{fig:convexidade-corda}. Partindo dessa geometria, chegaremos à
desigualdade de Jensen, que compara $g(\mathbb{E}[X])$ com
$\mathbb{E}[g(X)]$ quando $g$ é convexa. A comparação entre momentos de
ordens distintas, dada pela desigualdade de Lyapunov, será uma de suas
aplicações.

\subsection{Funções Convexas e Algumas de Suas Caracterizações}\label{subsec:convexidade-definicao}

\begin{definition}[Função convexa]\label{def:convexa}
Sejam $-\infty\leqslant a<b\leqslant\infty$. Uma função
$g\colon(a,b)\to\mathbb{R}$ é \emph{convexa} se
\begin{equation}\label{eq:convexidade}
g\bigl((1-t)u+tv\bigr)\leqslant(1-t)g(u)+tg(v)
\end{equation}
para todos $t\in[0,1]$ e $u,v\in(a,b)$. Se a desigualdade é estrita sempre que
$u\neq v$ e $t\in(0,1)$, dizemos que $g$ é \emph{estritamente convexa}.
\end{definition}

O argumento da função $g$ em \eqref{eq:convexidade} é ponto $(1-t)u+tv$.
Para verificar que o lado esquerdo de \eqref{eq:convexidade} está bem-definido,
basta observar que para qualquer valor de $t\in [0,1]$, temos que este ponto pertence
ao intervalo fechado $[u,v]$. Mais precisamente,
quando o parâmetro $t$ varia em $[0,1]$
este ponto percorre todo segmento de reta entre os pontos $u$ e $v$, veja figura abaixo

\begin{figure}[H]
\centering
\begin{tikzpicture}[x=1.5cm,y=1cm,>=latex,font=\small]
  % Um valor intermediário de t, com u<v.
  \draw[thin,->] (-0.7,0) -- (7.2,0) node[below] {$x$};
  \draw[thick] (0,0) -- (6,0);
  \draw[very thick,teal] (0,0) -- (2.2,0);
  \fill (0,0) circle (1.8pt);
  \fill (6,0) circle (1.8pt);
  \fill[teal] (2.2,0) circle (2.2pt);
  \node[below=5pt] at (0,0) {$u$};
  \node[below=5pt] at (6,0) {$v$};
  \node[below=5pt,teal] at (2.2,0) {$x_t=(1-t)u+tv$};

  % As distâncias explicam o papel do parâmetro t.
  \draw[decorate,decoration={brace,amplitude=5pt},teal]
    (0,0.18) -- (2.2,0.18)
    node[midway,above=7pt] {$t(v-u)$};
  \draw[decorate,decoration={brace,amplitude=5pt}]
    (2.2,0.18) -- (6,0.18)
    node[midway,above=7pt] {$(1-t)(v-u)$};
  \draw[->,teal] (0,-1.05) -- (6,-1.05)
    node[midway,below=3pt,font=\scriptsize]
      {$t$ cresce de $0$ a $1$};
\end{tikzpicture}
\caption{Para $u<v$, escrever $x_t=(1-t)u+tv=u+t(v-u)$ mostra que
$t$ é a fração do caminho de $u$ até $v$ já percorrida. Quando $t=0$,
temos $x_t=u$; quando $t=1$, temos $x_t=v$. Para $0<t<1$, o ponto
fica entre os extremos; em particular, $t=1/2$ fornece o ponto médio.}
\label{fig:combinacao-convexa}
\end{figure}

O lado direito de \eqref{eq:convexidade} é a altura do ponto
correspondente na corda que une $(u,g(u))$ e $(v,g(v))$. A definição diz,
portanto, que a corda fica acima do gráfico. A \Cref{fig:convexidade-corda}
mostra a situação e destaca a folga entre a corda e o gráfico.

% F1 do levantamento 05-convexidade: a corda fica acima do gráfico.
% g(x) = 0,4x^2 + 0,4 em [0,3]; u = 0,5, v = 2,7, x_0 = 1,6.
% g(u) = 0,500; g(v) = 3,316; g(x_0) = 1,424; corda em x_0 = 1,908.
\begin{figure}[H]
\centering
\begin{tikzpicture}[x=2.2cm, y=1.5cm, >=latex,
    rot/.style={font=\scriptsize},
    marca/.style={font=\scriptsize, text=muted}]

  % eixos
  \draw[thin, ->] (-0.15,0) -- (3.2,0) node[below, font=\scriptsize] {$x$};
  \draw[thin, ->] (0,-0.1) -- (0,4.8);

  % gráfico de g
  \draw[thick, domain=0:3, samples=120] plot (\x,{0.4*\x*\x+0.4});
  \node[rot, anchor=south west] at (3.02,4.0) {$y=g(x)$};

  % corda pelos pontos de abscissas u e v
  \draw[thick] (0.5,0.5) -- (2.7,3.316);
  \node[rot, rotate=41, anchor=south west] at (0.55,1.15) {corda};

  % extremos da corda
  \fill (0.5,0.5) circle (1.6pt);
  \node[rot, anchor=north west] at (0.55,0.44) {$g(u)$};
  \fill (2.7,3.316) circle (1.6pt);
  \node[rot, anchor=north west] at (2.74,3.2) {$g(v)$};

  % folga entre a corda e o gráfico no ponto de abscissa x_0
  \draw[thin, dashed, muted] (1.6,0) -- (1.6,1.424);
  \fill[teal] (1.6,1.424) circle (1.5pt);
  \fill[teal] (1.6,1.908) circle (1.5pt);
  \draw[very thick, teal] (1.6,1.424) -- (1.6,1.908);
  \draw[decorate, decoration={brace, amplitude=4pt, mirror}, teal]
      (1.7,1.43) -- (1.7,1.90);

  % marcas no eixo
  \foreach \p in {0.5,1.6,2.7} \draw[thin] (\p,0) -- (\p,-0.07);
  \node[marca, below] at (0.5,-0.05) {$u$};
  \node[marca, below] at (2.7,-0.05) {$v$};
  \node[marca, below] at (1.6,-0.05) {$x_0$};
\end{tikzpicture}
\caption{A função é convexa quando a corda que une dois pontos quaisquer do
gráfico fica acima do gráfico. O segmento vertical destacado mede a folga
$(1-t)g(u)+tg(v)-g(x_0)$, onde
$x_0=(1-t)u+tv$.}
\label{fig:convexidade-corda}
\end{figure}


\begin{remark}\label{rem:convexidade-lambda}
Também é comum escrever a combinação convexa de dois pontos na forma
$\lambda u+(1-\lambda)v$, com $\lambda\in[0,1]$, o que equivale a tomar
$\lambda=1-t$. No texto usaremos $(1-t)u+tv$ para dois pontos e reservaremos a
notação $\lambda_1,\ldots,\lambda_n$ para combinações convexas de $n$ pontos.
\end{remark}

A desigualdade \eqref{eq:convexidade} da \Cref{def:convexa}
se estende a qualquer número finito de pontos. Se
$g\colon(a,b)\to\mathbb{R}$ é uma função convexa, se $x_1,\ldots,x_n\in(a,b)$ e se $\lambda_1,\ldots,\lambda_n$
são números reais não negativos com soma $1$, então
\begin{equation}\label{eq:convexidade-n-pontos}
g\Bigl(\sum_{i=1}^{n}\lambda_ix_i\Bigr)
\leqslant
\sum_{i=1}^{n}\lambda_ig(x_i).
\end{equation}
A definição de convexidade compara apenas dois pontos. Para usá-la em uma
combinação de vários pontos, a ideia é agrupar alguns deles e tratar o grupo
como um único ponto. Por exemplo, podemos escrever
\[
\frac14x_1+\frac14x_2+\frac12x_3
=\frac12\underbrace{\left(\frac12x_1+\frac12x_2\right)}_{y}
+\frac12x_3.
\]
Primeiro formamos o ponto médio $y$ de $x_1$ e $x_2$; depois tomamos o ponto
médio de $y$ e $x_3$. Os pesos $1/4$ aparecem porque o peso $1/2$ do grupo
é dividido igualmente entre seus dois pontos. Assim, uma combinação de três
pontos é construída por duas combinações de dois pontos, como mostra a
\Cref{fig:convexidade-agrupamento}.

\begin{figure}[H]
\centering
\begin{tikzpicture}[x=1.4cm,y=1cm,>=latex,font=\small,
  ponto/.style={fill=white,inner sep=4pt},
  peso/.style={fill=white,inner sep=2pt,font=\scriptsize}]
  \node[ponto] (x1) at (0,2) {$x_1$};
  \node[ponto] (x2) at (3,2) {$x_2$};
  \node[ponto] (x3) at (7,2) {$x_3$};
  \node[ponto,text=teal] (y) at (1.5,0.5)
    {$y=\frac12x_1+\frac12x_2$};
  \node[ponto] (z) at (4,-1.2)
    {$\frac12y+\frac12x_3=\frac14x_1+\frac14x_2+\frac12x_3$};
  \draw[->,thick,teal] (x1) -- node[peso] {$\frac12$} (y);
  \draw[->,thick,teal] (x2) -- node[peso] {$\frac12$} (y);
  \draw[->,thick] (y) -- node[peso] {$\frac12$} (z);
  \draw[->,thick] (x3) -- node[peso] {$\frac12$} (z);
\end{tikzpicture}
\caption{Uma combinação de três pontos construída em duas etapas.
Os números nas setas são os pesos usados em cada etapa. Cada um dos
pontos $x_1$ e $x_2$ recebe metade do peso de $y$, ficando com peso
final $(1/2)(1/2)=1/4$. O ponto $x_3$ recebe diretamente peso $1/2$.
O diagrama representa o agrupamento dos pontos, sem pressupor uma ordem
entre seus valores na reta.}
\label{fig:convexidade-agrupamento}
\end{figure}


A convexidade pode acompanhar essas duas etapas. Aplicando a definição
primeiro ao par $y,x_3$ e depois ao par $x_1,x_2$, obtemos
\[
\begin{aligned}
g\left(\frac14x_1+\frac14x_2+\frac12x_3\right)
&\leqslant\frac12g(y)+\frac12g(x_3)\\[2pt]
&\leqslant\frac12\left(\frac12g(x_1)+\frac12g(x_2)\right)
  +\frac12g(x_3)\\[2pt]
&=\frac14g(x_1)+\frac14g(x_2)+\frac12g(x_3).
\end{aligned}
\]

Para pesos quaisquer, usaremos a mesma ideia em uma prova por indução.
O caso $n=2$ é a própria definição de convexidade. Suponha que a
desigualdade já seja válida para combinações de $n-1$ pontos e consideremos
uma combinação de $n$ pontos. Vamos agrupar os primeiros $n-1$ pontos,
deixando $x_n$ separado. O peso total do grupo é
\[
\lambda_1+\cdots+\lambda_{n-1}=1-\lambda_n.
\]
Se $\lambda_n=1$, todos os outros pesos são zero e os dois lados de
\eqref{eq:convexidade-n-pontos} são iguais a $g(x_n)$. Basta, portanto,
tratar o caso $\lambda_n<1$, no qual o peso do grupo é positivo.

Para representar esse grupo por uma combinação convexa, precisamos de
pesos cuja soma seja $1$. Os pesos originais somam $1-\lambda_n$;
dividindo cada um deles por essa soma, obtemos
\[
\alpha_i=\frac{\lambda_i}{1-\lambda_n},\qquad 1\leqslant i\leqslant n-1.
\]
Esses novos pesos são não negativos e satisfazem
\[
\sum_{i=1}^{n-1}\alpha_i
=\frac{\lambda_1+\cdots+\lambda_{n-1}}{1-\lambda_n}=1.
\]
Podemos, assim, definir o ponto que representa o grupo por
\[
y=\sum_{i=1}^{n-1}\alpha_ix_i.
\]
Como $y$ é uma média ponderada dos primeiros $n-1$ pontos, ele fica entre
o menor e o maior deles e, em particular, pertence a $(a,b)$. Ao atribuir
ao grupo seu peso total $1-\lambda_n$, recuperamos os pesos originais:
\[
(1-\lambda_n)y+\lambda_nx_n
=\sum_{i=1}^{n-1}(1-\lambda_n)\alpha_ix_i+\lambda_nx_n
=\sum_{i=1}^n\lambda_ix_i.
\]
Agora temos uma combinação dos dois pontos $y$ e $x_n$, à qual podemos
aplicar a definição de convexidade:
\[
g\left(\sum_{i=1}^n\lambda_ix_i\right)
\leqslant(1-\lambda_n)g(y)+\lambda_ng(x_n).
\]
Resta estimar $g(y)$. Como $y$ é uma combinação convexa de $n-1$ pontos,
a hipótese de indução fornece
\[
g(y)\leqslant\sum_{i=1}^{n-1}\alpha_i g(x_i).
\]
Substituindo essa estimativa na desigualdade anterior e usando
$(1-\lambda_n)\alpha_i=\lambda_i$, concluímos que
\[
g\left(\sum_{i=1}^n\lambda_ix_i\right)
\leqslant(1-\lambda_n)\sum_{i=1}^{n-1}\alpha_i g(x_i)+\lambda_ng(x_n)
=\sum_{i=1}^n\lambda_i g(x_i).
\]
Isso completa o passo de indução e demonstra
\eqref{eq:convexidade-n-pontos}.

Essa desigualdade já tem uma interpretação probabilística. Se $X$ assume
um número finito de valores $x_1,\ldots,x_n$ e
$p_i=\mathbb{P}(X=x_i)$, então os $p_i$ são não negativos e somam $1$.
Portanto, podemos usá-los como pesos em
\eqref{eq:convexidade-n-pontos}, obtendo
\[
g\bigl(\mathbb{E}[X]\bigr)
=g\left(\sum_{i=1}^n p_ix_i\right)
\leqslant\sum_{i=1}^n p_i g(x_i)
=\mathbb{E}[g(X)].
\]
Assim, para variáveis com um número finito de valores, a comparação entre
a função da média e a média da função decorre diretamente da convexidade.
Essa é a desigualdade de Jensen nesse caso particular. Para tratar também
variáveis com infinitos valores ou com densidade, desenvolveremos a
geometria da convexidade até obter uma reta de suporte. Ela permitirá
demonstrar Jensen sob as hipóteses de integrabilidade da
\Cref{subsec:jensen}, sem precisar reduzir cada distribuição a uma soma
finita. Começaremos estudando como variam as inclinações das secantes ao
gráfico de uma função convexa.

\subsection{As Inclinações das Retas Secantes}\label{subsec:secantes}

Até aqui, usamos a convexidade para comparar a altura do gráfico com a
altura de uma corda. Vamos agora examinar a inclinação das retas que passam
por dois pontos do gráfico, chamadas de \emph{retas secantes}. Se as
abscissas desses pontos são $u<v$, a inclinação da reta secante associada a estes pontos é
dada por
\[
\frac{g(v)-g(u)}{v-u}.
\]
O numerador mede a variação na altura e o denominador, o deslocamento
horizontal. Essa razão é também a taxa média de variação de $g$ entre
$u$ e $v$.

Para entender o que a convexidade diz sobre essa taxa, vamos considerar um ponto
intermediário $w$, com $u<w<v$. Como $(w,g(w))$ fica abaixo da corda entre
$(u,g(u))$ e $(v,g(v))$, a secante de $u$ a $w$ tem inclinação menor ou
igual à de $u$ a $v$. Pelo mesmo motivo, a secante de $w$ a $v$ tem
inclinação maior ou igual à de $u$ a $v$. O lema a seguir traduz essa
observação em desigualdades, ilustradas na \Cref{fig:convexidade-secantes}.

% F2 do levantamento 05-convexidade: as inclinações das secantes crescem.
% g(x) = 0,4x^2 + 0,4; u = 0,4, w = 1,3, v = 2,6.
% g(u) = 0,464; g(w) = 1,076; g(v) = 3,104.
% s_1 = 0,680; s_2 = 1,562; s_uv = 1,200.
\begin{figure}[htbp]
\centering
\begin{tikzpicture}[x=2.0cm, y=1.35cm, >=latex,
    rot/.style={font=\scriptsize},
    marca/.style={font=\scriptsize, text=muted}]

  % eixos
  \draw[thin, ->] (-0.15,0) -- (3.2,0) node[below, font=\scriptsize] {$x$};
  \draw[thin, ->] (0,-0.1) -- (0,4.6);

  % gráfico de g
  \draw[thick, domain=0:3, samples=120] plot (\x,{0.4*\x*\x+0.4});

  % secantes
  \draw[thick] (0.4,0.464) -- (1.3,1.076);
  \draw[thick] (1.3,1.076) -- (2.6,3.104);
  \draw[thick, dashed, muted] (0.4,0.464) -- (2.6,3.104);

  % pontos
  \fill (0.4,0.464) circle (1.5pt);
  \fill (1.3,1.076) circle (1.5pt);
  \fill (2.6,3.104) circle (1.5pt);

  % marcas no eixo
  \foreach \p in {0.4,1.3,2.6} \draw[thin] (\p,0) -- (\p,-0.07);
  \node[marca, below] at (0.4,-0.05) {$u$};
  \node[marca, below] at (1.3,-0.05) {$w$};
  \node[marca, below] at (2.6,-0.05) {$v$};

  % quadro com as inclinações, na região livre acima da curva
  \node[anchor=north west, align=left, font=\scriptsize,
        fill=white, draw=border, rounded corners=1pt,
        inner sep=4pt] at (0.12,4.45)
    {\textcolor{black}{\rule{0.7cm}{0.6pt}}\quad $s_1\approx 0{,}68$\\[2pt]
     \textcolor{black}{\rule{0.7cm}{0.6pt}}\quad $s_2\approx 1{,}56$\\[2pt]
     \textcolor{muted}{\rule{0.7cm}{0.6pt}}\quad $s_{uv}\approx 1{,}20$\\[3pt]
     \hspace*{0.1cm}$s_1\leqslant s_{uv}\leqslant s_2$};
\end{tikzpicture}
\caption{Para três pontos $u<w<v$, a inclinação $s_1$ da secante sobre $[u,w]$
não excede a inclinação $s_{uv}$ da secante sobre $[u,v]$, que por sua vez não
excede a inclinação $s_2$ da secante sobre $[w,v]$. As inclinações das secantes
crescem com o deslocamento do ponto para a direita.}
\label{fig:convexidade-secantes}
\end{figure}


Essa ordem entre as inclinações será útil
para construir uma reta que passe por $(w,g(w))$ e fique abaixo de todo
o gráfico: sua inclinação deverá ser pelo menos tão grande quanto as das
secantes à esquerda de $w$ e não maior que as das secantes à direita.

\begin{lemma}[Monotonicidade das secantes]\label{lem:secantes}
Seja $g\colon(a,b)\to\mathbb{R}$ convexa e sejam $u<w<v$ em $(a,b)$. Então
\begin{equation}\label{eq:secantes}
\frac{g(w)-g(u)}{w-u}
\leqslant
\frac{g(v)-g(u)}{v-u}
\leqslant
\frac{g(v)-g(w)}{v-w}.
\end{equation}
\end{lemma}

\begin{proof}
Primeiro, vamos representar o ponto $w$ como uma combinação convexa de $u$ e $v$, isto é, $w=(1-t)u+tv$, onde
\[
t=\frac{w-u}{v-u}\in(0,1).
\]
Pela convexidade,
\begin{equation}\label{eq:secantes-convexidade}
g(w)\leqslant(1-t)g(u)+tg(v).
\end{equation}
Subtraindo $g(u)$ dos dois lados de \eqref{eq:secantes-convexidade}, obtemos
\[
g(w)-g(u)
\leqslant
t\bigl(g(v)-g(u)\bigr),
\]
e, dividindo por $w-u>0$ e usando $t=(w-u)/(v-u)$, segue que
\[
\frac{g(w)-g(u)}{w-u}
\leqslant
\frac{g(v)-g(u)}{v-u}.
\]
Isso prova a primeira desigualdade. 


Para a segunda, queremos comparar as
variações de $g$ entre $u$ e $v$ e entre $w$ e $v$. Por isso, subtraímos
$g(v)$ dos dois lados de \eqref{eq:secantes-convexidade}, obtendo
\[
g(w)-g(v)
\leqslant
(1-t)\bigl(g(u)-g(v)\bigr).
\]
Multiplicando por $(-1)$, ambos os lados da desigualdade acima ficamos com 
\[
g(v)-g(w)
\geqslant
(1-t)\bigl(g(v)-g(u)\bigr)
=\frac{v-w}{v-u}\bigl(g(v)-g(u)\bigr),
\]
onde na última igualdade, usamos
\[
1-t=1-\frac{w-u}{v-u}=\frac{v-w}{v-u}.
\]
Assim, ao dividir por $v-w>0$, o fator $v-w$ do lado direito se cancela,
e obtemos a segunda desigualdade:
\[
\frac{g(v)-g(u)}{v-u}
\leqslant
\frac{g(v)-g(w)}{v-w}.
\]
\end{proof}

Para $w\in(a,b)$ fixado vamos ver como variam as inclinações das secantes que
têm uma extremidade no ponto $(w,g(w))$. À esquerda de $w$, isto é, 
ao longo do intervalo $(a,w)$ a função
\[
u\longmapsto\frac{g(w)-g(u)}{w-u}
\]
é não-decrescente. De fato, se $a<u_1<u_2<w$, a segunda
desigualdade de \eqref{eq:secantes}, aplicada ao trio $u_1<u_2<w$, dá
\[
\frac{g(w)-g(u_1)}{w-u_1}
\leqslant
\frac{g(w)-g(u_2)}{w-u_2}.
\]
Similarmente, à direita de $w$, isto é, no intervalo $(w,b)$ a função
\[
v\longmapsto\frac{g(v)-g(w)}{v-w}
\]
é não-decrescente. De fato, se $w<v_1<v_2<b$, a primeira
desigualdade de \eqref{eq:secantes}, aplicada ao trio $w<v_1<v_2$, fornece
\[
\frac{g(v_1)-g(w)}{v_1-w}
\leqslant
\frac{g(v_2)-g(w)}{v_2-w}.
\]

Agora precisamos verificar que essas inclinações não divergem quando
as variáveis se aproximam de $w$. Para isso, usamos outra consequência
de \eqref{eq:secantes}: toda secante à esquerda tem inclinação menor ou
igual à de qualquer secante à direita, isto é,
\begin{equation}\label{eq:secantes-esquerda-direita}
\frac{g(w)-g(u)}{w-u}
\leqslant
\frac{g(v)-g(w)}{v-w}
\qquad (a<u<w<v<b).
\end{equation}
Fixando $v=v_0\in(w,b)$, obtemos uma cota superior finita para todas as
inclinações à esquerda. Fixando $u=u_0\in(a,w)$, obtemos uma cota inferior
finita para todas as inclinações à direita. Portanto, os números
\[
L_w=\sup_{u\in(a,w)}\frac{g(w)-g(u)}{w-u},
\qquad
R_w=\inf_{v\in(w,b)}\frac{g(v)-g(w)}{v-w}
\]
são finitos. Pela monotonicidade que acabamos de provar, eles são também
os limites laterais
\[
L_w=\lim_{u\to w^-}\frac{g(w)-g(u)}{w-u},
\qquad
R_w=\lim_{v\to w^+}\frac{g(v)-g(w)}{v-w}.
\]
À esquerda, aproximar $u$ de $w$ faz as inclinações crescerem até seu
supremo; à direita, aproximar $v$ de $w$ faz as inclinações decrescerem até
seu ínfimo.

Além disso, tomando primeiro o supremo com $u\in (a,w)$ em \eqref{eq:secantes-esquerda-direita}, obtemos 
\begin{align*}
L_{w}\equiv 
\sup_{u\in(a,w)}\frac{g(w)-g(u)}{w-u}
\leqslant
\sup_{u\in(a,w)}
\frac{g(v)-g(w)}{v-w}
=
\frac{g(v)-g(w)}{v-w},
\end{align*}
para todo $v\in (w,b)$. 
Tomando agora, na desigualdade acima, o ínfimo com $v\in (w,b)$ concluímos que
\begin{align*}
L_{w}
=
\inf_{v\in(w,b)} L_{w} 
\leqslant
\inf_{v\in(w,b)} \frac{g(v)-g(w)}{v-w}
\equiv 
R_{w}.
\end{align*}

Assim, existe pelo menos um número $\alpha$ que é maior ou igual a todas
as inclinações à esquerda e menor ou igual a todas as inclinações à
direita. Na próxima subseção, veremos que os valores
$\alpha\in[L_w,R_w]$ são precisamente as inclinações das retas que passam
por $(w,g(w))$ e ficam abaixo de todo o gráfico de $g$.

\subsection{A Reta Suporte}\label{subsec:reta-suporte}

Seja $g\colon(a,b)\to\mathbb{R}$ uma função convexa. Uma \emph{reta de
suporte} ao seu gráfico em $w\in(a,b)$ é uma reta que passa por
$(w,g(w))$ e cuja altura, em cada abscissa do domínio, é menor ou igual
à altura do gráfico. O próximo teorema mostra que a convexidade garante
a existência de pelo menos uma reta assim em cada ponto do gráfico.
Se além de ser uma função convexa $g$ é também diferenciável em $w$, então 
existe uma única reta suporte nesse ponto que é a reta tangente ao gráfico neste ponto. 
É importante observar que convexidade garante a 
existência de uma reta suporte mesmo quando a derivada não existe. No caso de não diferenciabilidade 
podemos ter a existência de várias retas suporte com inclinações distintas.

\begin{theorem}[Reta suporte]\label{thm:reta-suporte}
Seja $g\colon(a,b)\to\mathbb{R}$ convexa e seja $w\in(a,b)$. Então existe
$\alpha\in\mathbb{R}$ tal que
\begin{equation}\label{eq:reta-suporte}
g(w)+\alpha(x-w) \leqslant g(x)
\qquad
\text{para todo }x\in(a,b).
\end{equation}
\end{theorem}

\begin{proof}
Sejam $L_w$ e $R_w$ como acima. Como $L_w\leqslant R_w$, escolha
$\alpha\in[L_w,R_w]$. Se $x<w$, então
\[
\frac{g(w)-g(x)}{w-x}\leqslant L_w\leqslant\alpha,
\]
e multiplicando por $w-x>0$ obtemos $g(w)-g(x)\leqslant\alpha(w-x)$, que implica
na desigualdade \eqref{eq:reta-suporte} para $x<w$. Se $x>w$, então
segue da definição de $R_w$ que 
\[
\alpha\leqslant R_w\leqslant\frac{g(x)-g(w)}{x-w},
\]
e multiplicando por $x-w>0$ obtemos \eqref{eq:reta-suporte} para $x>w$. Para
$x=w$ a desigualdade é uma igualdade.
\end{proof}

A \Cref{fig:convexidade-suporte} mostra dois casos. À esquerda, a função é
diferenciável e a reta de suporte é a tangente, com $\alpha=g'(w)$. À direita,
a função tem um bico e existem várias retas de suporte, uma para cada
$\alpha\in[L_w,R_w]$. É por isso que a convexidade não exige derivada.

% F3 do levantamento 05-convexidade: a reta de suporte, com e sem derivada.
% Painel esquerdo: g(x) = 0,4x^2 + 0,4; w = 1,4; g'(w) = 1,12.
% Painel direito: g(x) = |x - 1,5| + 0,7, com bico em w = 1,5.
\begin{figure}[htbp]
\centering
\begin{minipage}{0.47\linewidth}
\centering
\begin{tikzpicture}[x=1.5cm, y=1.1cm, >=latex,
    rot/.style={font=\scriptsize},
    marca/.style={font=\scriptsize, text=muted}]

  \draw[thin, ->] (-0.15,0) -- (3.2,0) node[below, font=\scriptsize] {$x$};
  \draw[thin, ->] (0,-0.6) -- (0,4.95);

  \draw[thick, domain=0:3, samples=120] plot (\x,{0.4*\x*\x+0.4});
  \node[rot, anchor=south west] at (3.02,4.0) {$y=g(x)$};

  % reta tangente em w
  \draw[teal, thick, domain=0.3:2.9] plot (\x,{1.12*\x-0.384});
  \fill (1.4,1.184) circle (1.6pt);
  \draw[thin, dashed, muted] (1.4,0) -- (1.4,1.184);
  \node[marca, below] at (1.4,-0.05) {$w$};

  % rótulo da reta de suporte, abaixo da tangente
  \node[rot, text=teal, rotate=39, anchor=north west] at (1.9,1.72)
    {reta de suporte};
\end{tikzpicture}
\end{minipage}\hfill
\begin{minipage}{0.47\linewidth}
\centering
\begin{tikzpicture}[x=1.5cm, y=1.1cm, >=latex,
    rot/.style={font=\scriptsize},
    marca/.style={font=\scriptsize, text=muted}]

  \draw[thin, ->] (0.1,0) -- (2.95,0) node[below, font=\scriptsize] {$x$};
  \draw[thin, ->] (1.5,0) -- (1.5,2.35);

  % função com bico
  \draw[thick] (0.2,2.1) -- (1.5,0.7) -- (2.8,2.1);

  % reta de suporte horizontal, que não é nenhum dos ramos
  \draw[teal, thick, dashed] (0.2,0.7) -- (2.8,0.7);
  \fill (1.5,0.7) circle (1.6pt);
  \node[marca, below] at (1.5,-0.05) {$w$};

  \node[rot, text=teal, anchor=south] at (2.55,0.75) {$\alpha=0$};
  \node[rot, text=muted, anchor=south] at (0.65,1.85) {$\alpha=-1$};
  \node[rot, text=muted, anchor=south] at (2.35,1.85) {$\alpha=+1$};
  \node[rot, anchor=north] at (1.5,0.55) {$\alpha\in[L_w,R_w]$};
\end{tikzpicture}
\end{minipage}
\caption{Toda função convexa admite em cada ponto uma reta de suporte. À
esquerda, o caso diferenciável, em que a reta de suporte é a tangente, de
inclinação $\alpha=g'(w)$. À direita, uma função com bico: qualquer reta de
inclinação $\alpha\in[L_w,R_w]$ que passa pelo bico é uma reta de suporte, e
por isso a convexidade não exige derivada.}
\label{fig:convexidade-suporte}
\end{figure}


\begin{remark}\label{rem:suporte-derivavel}
Se $g$ é diferenciável em $w$, então $L_w=R_w=g'(w)$ e a reta de suporte é
única. De fato, o quociente $(g(w)-g(u))/(w-u)$ tende a $g'(w)$ quando
$u\to w^-$ e $(g(v)-g(w))/(v-w)$ tende a $g'(w)$ quando $v\to w^+$. Como as
duas funções são monótonas e limitadas, os supremo e ínfimo coincidem com os
limites laterais.
\end{remark}

\subsection{Caracterizações Equivalentes de Convexidade}\label{subsec:caracterizacoes}

Reunimos agora as três descrições equivalentes do conceito de convexidade e acrescentamos 
a caracterização via derivada segunda para o caso de funções convexas suaves
que aparecem frequentemente nos exemplos.

\begin{proposition}\label{prop:caracterizacoes}
Seja $g\colon(a,b)\to\mathbb{R}$. As afirmações seguintes são equivalentes.
\begin{enumerate}[label=(\roman*)]
\item $g$ é convexa.
\item Para todos $u<w<v$ em $(a,b)$, vale \eqref{eq:secantes}.
\item Para cada $w\in(a,b)$ existe $\alpha\in\mathbb{R}$ tal que vale
\eqref{eq:reta-suporte}.
\end{enumerate}
Se $g$ é duas vezes continuamente diferenciável, as três são ainda equivalentes a
\begin{enumerate}[label=(\roman*),start=4]
\item $g''(x)\geqslant0$ para todo $x\in(a,b)$.
\end{enumerate}
\end{proposition}

\begin{proof}
A implicação (i)$\Rightarrow$(ii) é o \Cref{lem:secantes}, e
(ii)$\Rightarrow$(iii) é o \Cref{thm:reta-suporte}, cuja prova usa apenas a
existência de $L_w$ e $R_w$ e a desigualdade $L_w\leqslant R_w$, que por sua
vez segue de (ii).

Para provar (iii)$\Rightarrow$(i), fixe $u,v\in(a,b)$ e $t\in[0,1]$ e
considere $w=(1-t)u+tv$. A ideia é usar a reta suporte em $w$ para comparar
$g(w)$ com a combinação convexa $(1-t)g(u)+tg(v)$. Por (iii), existe
$\alpha\in\mathbb{R}$ tal que
\[
g(u)\geqslant g(w)+\alpha(u-w)
\quad\text{e}\quad
g(v)\geqslant g(w)+\alpha(v-w).
\]
Multiplicando essas desigualdades pelos escalares não-negativos $1-t$ e $t$,
respectivamente, e somando, obtemos
\[
(1-t)g(u)+tg(v)
\geqslant
g(w)+\alpha\bigl((1-t)(u-w)+t(v-w)\bigr)
=
g(w),
\]
pois $(1-t)(u-w)+t(v-w)=(1-t)u+tv-w=0$, pela escolha de $w$.
Como $w=(1-t)u+tv$, essa é precisamente a desigualdade de convexidade
\eqref{eq:convexidade}.

\bigskip 
Suponha agora que $g$ seja duas vezes continuamente diferenciável. Como
(i), (ii) e (iii) já são equivalentes, basta provar as implicações
(iv)$\Rightarrow$(ii) e (i)$\Rightarrow$(iv).

Para provar (iv)$\Rightarrow$(ii), suponha que $g''\geqslant0$. Então
$g'$ é não decrescente. Fixe $u<w<v$ em $(a,b)$. Pelo teorema do valor médio,
a inclinação da secante
sobre $[u,w]$ é $g'(\xi_1)$ para algum $\xi_1\in(u,w)$ e a inclinação sobre
$[w,v]$ é $g'(\xi_2)$ para algum $\xi_2\in(w,v)$. Como $\xi_1<\xi_2$ e $g'$ é
não decrescente, vale \eqref{eq:secantes}, o que prova (ii).

Para provar (i)$\Rightarrow$(iv), suponha que $g$ seja convexa.
Seja $x_0\in(a,b)$ e suponha, por
absurdo, que $g''(x_0)<0$. Pela continuidade de $g''$, existe $h>0$ tal que
$g''(\xi)<0$, para todo $\xi\in (x_0-h,x_0+h)$. A fórmula de Taylor com resto de Lagrange, aplicada
a $x_0+h$ e a $x_0-h$, dá
\[
g(x_0+h)+g(x_0-h)-2g(x_0)
=
\frac{h^2}{2}\bigl(g''(\xi_+)+g''(\xi_-)\bigr)
\]
para certos $\xi_+\in(x_0,x_0+h)$ e $\xi_-\in(x_0-h,x_0)$. O lado direito é
negativo, logo
\[
g(x_0)>\frac{g(x_0+h)+g(x_0-h)}{2},
\]
o que contradiz a convexidade no par $x_0-h<x_0+h$ com $t=1/2$.
\end{proof}

\begin{remark}\label{rem:convexidade-exemplos}
A proposição acima fornece um critério rápido para verificar convexidade. 
Por exemplo, como consequência deste critério, verificamos imediatamente que são convexas as funções
$g(x)= x^2$, $g(x)= e^{tx}$, para cada $t\in\mathbb{R}$ fixado, $g(x)= x^p$, para cada $p\geqslant1$ fixado, 
$g(x)=1/x$ em $(0,\infty)$ e $g(x)=-\log x$ em
$(0,\infty)$. A última tem derivada segunda $1/x^2>0$ e este é o exemplo que
produz a desigualdade entre as médias geométrica e aritmética.
\end{remark}

\subsection{A Desigualdade de Jensen}\label{subsec:jensen}

Usando o resultado de existência da reta suporte, 
a prova da Desigualdade de Jensen é curta. Basta aplicar a desigualdade
\eqref{eq:reta-suporte} ponto a ponto e tomar esperanças.

\begin{theorem}[Desigualdade de Jensen]\label{thm:jensen}
Seja $X$ uma variável aleatória assumindo valores em um intervalo
$(a,b)$, com $\mathbb{E}[X]$ finita, e seja $g\colon(a,b)\to\mathbb{R}$
convexa tal que $\mathbb{E}\bigl[|g(X)|\bigr]<\infty$. Então
\begin{equation}\label{eq:jensen}
\mathbb{E}\bigl[g(X)\bigr]\geqslant g\bigl(\mathbb{E}[X]\bigr).
\end{equation}
\end{theorem}

\begin{proof}
Seja $\mu=\mathbb{E}[X]$. Como $a< X< b$ temos da monotonicidade da esperança 
que $a< \mathbb{E}[X] < b$. Logo $\mu\in (a,b)$. Aplicando o 
\Cref{thm:reta-suporte}, com $w=\mu$, podemos afirmar que existe $\alpha\in\mathbb{R}$ tal
que
\[
g(x)\geqslant g(\mu)+\alpha(x-\mu)
\qquad
\text{para todo }x\in(a,b).
\]
Substituindo $x$ por $X$ e tomando esperanças, obtemos
\[
\mathbb{E}\bigl[g(X)\bigr]
\geqslant
g(\mu)+\alpha\bigl(\mathbb{E}[X]-\mu\bigr)
=
g(\mu) = g\bigl(\mathbb{E}[X]\bigr),
\]
o que prova a validade da desigualdade \eqref{eq:jensen}.
\end{proof}

A mesma demonstração se estende a qualquer intervalo $I\subseteq\mathbb{R}$
(aberto, fechado ou semi-aberto) no qual $g$ seja convexa e contínua nos
extremos pertencentes a $I$: se $\mu=\mathbb{E}[X]$ está no interior de $I$, a
desigualdade \eqref{eq:reta-suporte} vale no interior e se estende aos
extremos por continuidade; se $\mu$ coincide com um extremo de $I$, então
$X=\mu$ quase certamente e \eqref{eq:jensen} vale com igualdade.

A \Cref{fig:jensen-esperanca} interpreta a desigualdade em termos de um
centro de massa. Suponha que $X$ assuma os valores $u$ e $v$, onde $u<v$, 
com probabilidades $p$ e $1-p$, respectivamente, onde $0<p<1$. Atribuindo essas massas aos
pontos $P_1=(u,g(u))$ e $P_2=(v,g(v))$, o centro de massa é
o ponto em $\mathbb{R}^2$ dado pela expressão
\[
A=pP_1+(1-p)P_2
 =\bigl(\mathbb{E}[X],\mathbb{E}[g(X)]\bigr).
\]
Esse ponto pertence à corda que une $P_1$ a $P_2$. Já o ponto
$B=\bigl(\mathbb{E}[X],g(\mathbb{E}[X])\bigr)$ pertence ao gráfico de $g$
e tem a mesma abscissa que $A$. Pela convexidade, $B$ está abaixo de $A$
ou coincide com ele. Assim, a desigualdade de Jensen afirma que a
diferença
\[
\Delta=\mathbb{E}[g(X)]-g(\mathbb{E}[X])
\]
é não negativa. Na figura abaixo, $\Delta$ é exatamente o comprimento do segmento de reta vertical
que une os ponto $B$ e $A$.

% Jensen como centro de massa: g(x) = 0,5x^2 + 0,35.
% u = 0,5; v = 2,7; p = 1/2; mu = 1,6.
% g(u) = 0,475; g(v) = 3,995; g(mu) = 1,63; E[g(X)] = 2,235.
\begin{figure}[H]
\centering
\begin{tikzpicture}[x=2.7cm, y=1.45cm, >=latex,
    rot/.style={font=\small},
    projecao/.style={thin, densely dashed, draw=muted!65}]

  \coordinate (P1) at (0.5,0.475);
  \coordinate (P2) at (2.7,3.995);
  \coordinate (A) at (1.6,2.235);
  \coordinate (B) at (1.6,1.63);

  \draw[thin, ->] (-0.12,0) -- (3.25,0) node[below, rot] {$x$};
  \draw[thin, ->] (0,-0.12) -- (0,4.9) node[left, rot] {$y$};

  % As projeções identificam as coordenadas dos dois pontos comparados.
  \draw[projecao] (0,2.235) node[left, rot] {$\mathbb{E}[g(X)]$} -- (A);
  \draw[projecao] (0,1.63) node[left, rot] {$g(\mathbb{E}[X])$} -- (B);
  \draw[projecao] (1.6,0) -- (B);
  \draw[projecao] (0.5,0) -- (P1);
  \draw[projecao] (2.7,0) -- (P2);
  \node[below, rot] at (0.5,0) {$u$};
  \node[below, rot] at (1.6,0) {$\mu=\mathbb{E}[X]$};
  \node[below, rot] at (2.7,0) {$v$};

  \draw[thick, domain=0:2.98, samples=100]
    plot (\x,{0.5*\x*\x+0.35});
  \node[rot, anchor=west] at (2.99,4.65) {$y=g(x)$};

  \draw[thick, muted] (P1) -- (P2)
    node[pos=0.73, above left, rot, text=muted] {corda};
  \fill (P1) circle (2pt);
  \fill (P2) circle (2pt);
  \node[rot, anchor=south east, align=right] at (0.5,0.5)
    {$P_1$};
  \node[rot, anchor=north west, align=left] at (2.77,3.97)
    {$P_2$};

  % A está na corda; B está no gráfico, na mesma abscissa.
  \draw[line width=1.8pt, teal] (B) -- (A);
  \fill[teal] (A) circle (2.5pt);
  \draw[thick, teal, fill=white] (B) circle (2.3pt);
  \node[rot, text=teal, above left=3pt] at (A) {$A$};
  \node[rot, text=teal, below right=3pt] at (B) {$B$};
  \node[rot, text=teal] at (1.74,2.08) {$\Delta$};
\end{tikzpicture}
\caption{Interpretação geométrica de Jensen. O centro de massa $A$ das
massas $p$ e $1-p$ está na corda; o ponto $B$ está no gráfico e tem a mesma
abscissa $\mu=\mathbb{E}[X]$. O segmento destacado tem comprimento
$\Delta=\mathbb{E}[g(X)]-g(\mathbb{E}[X])\geqslant0$.}
\label{fig:jensen-esperanca}
\end{figure}


\begin{remark}\label{rem:jensen-media}
Há duas maneiras de combinar a função $g$ com a esperança: podemos aplicar
$g$ aos valores de $X$ e depois tomar a média, obtendo $\mathbb{E}[g(X)]$,
ou tomar primeiro a média de $X$ e depois aplicar $g$, obtendo
$g(\mathbb{E}[X])$. Quando $g$ é convexa, a desigualdade de Jensen garante
que o primeiro resultado é maior ou igual ao segundo.

A diferença entre duas expressões que aparecem na Desigualdade de Jensen quantifica, 
assim, o efeito de substituir $X$ por
sua média antes de aplicar $g$. Esse efeito depende tanto da distribuição
de $X$ quanto da função escolhida. Por exemplo, para $g(x)=x^2$, a diferença
é exatamente a variância de $X$, quando o segundo momento é finito. Já para
uma função afim, ela é sempre nula, mesmo que $X$ não seja constante.
\end{remark}

\subsection{O Caso de Igualdade}\label{subsec:jensen-igualdade}

Para descobrir quando a desigualdade de Jensen é saturada, isto é, 
quando ocorre a igualdade, comparamos $g(X)$ com a
reta de suporte ao gráfico de $g$ no ponto $\mu=\mathbb{E}[X]$. A diferença
entre esses valores é sempre não negativa, e sua esperança é exatamente a
diferença entre os dois lados de \eqref{eq:jensen}. Se $g$ é estritamente
convexa, essa reta toca o gráfico apenas em $\mu$; assim, a igualdade só pode
ocorrer quando $X=\mu$ quase certamente.

\begin{proposition}\label{prop:jensen-igualdade}
Nas hipóteses do \Cref{thm:jensen}, suponha que $g$ seja estritamente convexa.
Então vale a igualdade em \eqref{eq:jensen} se, e somente se,
$X=\mathbb{E}[X]$ quase certamente.
\end{proposition}

\begin{proof}
Se $X=\mu$ quase certamente, então $\mathbb{E}[g(X)]=g(\mu) = g\bigl(\mathbb{E}[X] \bigr)$ e a igualdade
vale. Suponha agora que valha a igualdade. Seja $\alpha$ o coeficiente dado
pela reta de suporte em $\mu$ e defina
\[
Z=g(X)-g(\mu)-\alpha(X-\mu).
\]
Por \eqref{eq:reta-suporte}, $Z\geqslant0$. Como estamos assumindo que 
vale a igualdade $\mathbb{E}[g(X)]= g\bigl(\mathbb{E}[X] \bigr)$
então temos que $\mathbb{E}[Z]=0$. Como mostrado anteriormente 
sabemos que se uma variável aleatória é não-negativa e possui esperança
nula então esta variável aleatória é nula quase certamente, logo
\[
g(X)=g(\mu)+\alpha(X-\mu)
\qquad
\text{quase certamente}.
\]
No evento $\{X\neq\mu\}$, a igualdade acima diria que a reta de suporte toca o
gráfico em dois pontos distintos, $X$ e $\mu$. Mas, pela convexidade estrita, o
ponto médio entre eles estaria estritamente abaixo da corda, e a corda coincide
com a reta de suporte, o que contradiz \eqref{eq:reta-suporte}. Logo
$\{X\neq\mu\}$ tem probabilidade nula.
\end{proof}

\subsection{Exemplos}\label{subsec:jensen-exemplos}

O primeiro exemplo recupera a não negatividade da variância.

\begin{example}[Não negatividade da variância]\label{ex:jensen-quadrado}
Tome $g(x)=x^2$, que é estritamente convexa. O \Cref{thm:jensen} dá
\[
\mathbb{E}[X^2]\geqslant\bigl(\mathbb{E}[X]\bigr)^2,
\]
isto é, $\operatorname{Var}(X)\geqslant0$. A igualdade vale se, e somente se,
$X$ é constante quase certamente. Esse é o mesmo conteúdo da desigualdade de
Cauchy-Schwarz para $X$ e $1$, que já foi provada nas notas sobre momentos, por outro
argumento. Aqui ela aparece como um caso particular de convexidade.
\end{example}

O segundo exemplo é a desigualdade entre as médias geométrica e aritmética.

\begin{example}[Média geométrica e média aritmética]\label{ex:media-geometrica}
A função $g(x)=-\log x$ é estritamente convexa em $(0,\infty)$, pois
$g''(x)=1/x^2>0$. Seja $X$ uma variável aleatória positiva com
$\mathbb{E}[X]<\infty$ e $\mathbb{E}\bigl[|\log X|\bigr]<\infty$. 
Então, segue da Desigualdade de Jensen que
\begin{equation}\label{eq:jensen-log}
\mathbb{E}[\log X]\leqslant\log\mathbb{E}[X],
\end{equation}
ou, exponenciando,
\[
\exp\bigl(\mathbb{E}[\log X]\bigr)\leqslant\mathbb{E}[X].
\]
Suponha agora que $X$ assume os valores positivos $x_1,\ldots,x_n$ com a mesma
probabilidade. Então
\[
\mathbb{E}[\log X]=\frac{1}{n}\sum_{i=1}^{n}\log x_i=\log\gamma,
\qquad
\mathbb{E}[X]=\frac{1}{n}\sum_{i=1}^{n}x_i=\bar{x},
\]
onde
\[
\gamma=\Bigl(\prod_{i=1}^{n}x_i\Bigr)^{1/n},
\qquad
\bar{x}=\frac{1}{n}\sum_{i=1}^{n}x_i
\]
são a média geométrica e a média aritmética dos $x_i$. Por \eqref{eq:jensen-log},
$\log\gamma\leqslant\log\bar{x}$, e como o logaritmo é crescente,
\[
\gamma\leqslant\bar{x}.
\]
Além disso, a igualdade vale se, e somente se, todos os $x_i$ são iguais.
\end{example}

O terceiro exemplo compara a média harmônica com a geométrica.

\begin{example}[Média harmônica]\label{ex:media-harmonica}
A média harmônica de números não-nulos $x_1,\ldots,x_n$ é definida por
\[
\frac{1}{\eta}=\frac{1}{n}\sum_{i=1}^{n}\frac{1}{x_i}.
\]
No caso em que $x_i>0$ podemos aplicar o \Cref{ex:media-geometrica} aos números $1/x_1,\ldots,1/x_n$,
obtemos
\[
\Bigl(\prod_{i=1}^{n}\frac{1}{x_i}\Bigr)^{1/n}
\leqslant
\frac{1}{n}\sum_{i=1}^{n}\frac{1}{x_i},
\]
isto é, $1/\gamma\leqslant1/\eta$, ou seja, $\eta\leqslant\gamma$. A média
harmônica não excede a geométrica, que por sua vez não excede a aritmética.
\end{example}

O quarto exemplo conecta Jensen com a função geradora de momentos.

\begin{example}[A função geradora de momentos]\label{ex:jensen-fgm}
Para cada $t\in\mathbb{R}\setminus\{0\}$, a função $g(x)=e^{tx}$ é
estritamente convexa (para $t=0$ é constante). Portanto segue da Desigualdade de Jensen que
\[
\mathbb{E}\bigl[e^{tX}\bigr]\geqslant e^{t\mathbb{E}[X]},
\]
isto é, $M_X(t)\geqslant e^{t\mathbb{E}[X]}$ para todo $t$ tal que a geradora
exista, com igualdade para $t\neq0$ se, e somente se, $X$ é constante quase
certamente.
\end{example}

\section{A Desigualdade de Lyapunov}\label{sec:lyapunov}

A convexidade também organiza a comparação entre momentos de ordens distintas.
A pergunta é a seguinte: se $\mathbb{E}|X|^r$ é finito para algum $r>0$, o que
isso diz sobre $\mathbb{E}|X|^s$ para outros valores de $s$? A resposta tem duas
partes. Primeiro, a finitude se propaga para baixo. Depois, as normas
$(\mathbb{E}|X|^r)^{1/r}$ crescem com $r$. A segunda parte é a desigualdade de
Lyapunov.

\subsection{Finitude para Ordens Menores}\label{subsec:lyapunov-finitude}

O ponto de partida é uma desigualdade pontual entre potências.

\begin{lemma}\label{lem:potencias}
Sejam $0<s\leqslant r$ e $x\in\mathbb{R}$. Então
\begin{equation}\label{eq:potencias}
|x|^s\leqslant1+|x|^r.
\end{equation}
\end{lemma}

\begin{proof}
Se $|x|\leqslant1$, então $|x|^s\leqslant1\leqslant1+|x|^r$. Se $|x|>1$,
então $|x|^s\leqslant|x|^r\leqslant1+|x|^r$, porque $s\leqslant r$ e a função
$p\longmapsto |x|^p$ é crescente em $(0,\infty)$ quando $|x|>1$.
\end{proof}

Tomando esperanças em \eqref{eq:potencias}, obtemos
\[
\mathbb{E}|X|^s\leqslant1+\mathbb{E}|X|^r.
\]
Em particular, se $\mathbb{E}|X|^r<\infty$, então $\mathbb{E}|X|^s<\infty$.
Isso mostra que o conjunto
\[
I=\bigl\{r>0:\mathbb{E}|X|^r<\infty\bigr\}
\]
é um intervalo com extremo esquerdo $0$: se $r\in I$, então $(0,r]\subseteq I$.
O intervalo pode ser $(0,\infty)$, como no caso gaussiano, em que todos os
momentos são finitos. Pode ser $(0,R)$ para algum $R>0$, como no caso da
distribuição de Cauchy, em que $\mathbb{E}|X|^r$ é finito para $0<r<1$ e
infinito para $r\geqslant1$. E pode ser $(0,R]$, se o momento de ordem $R$ for
finito mas os de ordem maior não forem. Em qualquer caso, a existência de um
momento de ordem alta implica a existência de todos os momentos de ordem menor.

\subsection{Enunciado e Prova da Desigualdade de Lyapunov}\label{subsec:lyapunov}

A desigualdade de Lyapunov compara as normas em ordens diferentes. A prova é
uma aplicação de Jensen à variável $Z^s$ e à função $y\longmapsto y^{r/s}$, que é
convexa porque $r/s\geqslant1$.

\begin{theorem}[Desigualdade de Lyapunov]\label{thm:lyapunov}
Seja $Z$ uma variável aleatória não negativa e sejam $0<s\leqslant r$. Então
\begin{equation}\label{eq:lyapunov}
\bigl(\mathbb{E}[Z^s]\bigr)^{1/s}
\leqslant 
\bigl(\mathbb{E}[Z^r]\bigr)^{1/r}
\end{equation}
\end{theorem}

\begin{proof}
Se $\mathbb{E}[Z^r]=\infty$, a desigualdade é imediata. Suponha que
$\mathbb{E}[Z^r]<\infty$. Pelo \Cref{lem:potencias}, aplicado a $Z$,
$\mathbb{E}[Z^s]<\infty$. Se $\mathbb{E}[Z^s]=0$, então $Z=0$ quase certamente
e a desigualdade é trivial. Suponha $\mathbb{E}[Z^s]>0$. Seja $Y=Z^s$ e
$g(y)=y^{r/s}$, definida para
$y\geqslant0$. Como $r/s\geqslant1$, a função $g$ é convexa em $[0,\infty)$.
De fato,
\begin{align*}
g'(y) = \frac{r}{s}\, y^{ \frac{r}{s} - 1 }
\qquad \text{e} \qquad 
g''(y) = \frac{r}{s} \left( \frac{r}{s}-1 \right) y^{ \frac{r}{s} - 2 }>0. 
\end{align*}

Usando a Desigualdade de Jensen para $Y$ e $g$ concluímos que
\[
\mathbb{E}\bigl[Z^r\bigr]
=
\mathbb{E}\bigl[Y^{r/s}\bigr]
=
\mathbb{E}\bigl[g(Y)\bigr]
\geqslant
g\bigl(\mathbb{E}[Y]\bigr)
=
\bigl(\mathbb{E}[Y]\bigr)^{r/s}
=
\bigl(\mathbb{E}[Z^s]\bigr)^{r/s}.
\]
Extraindo a raiz $r$-ésima dos dois lados, que é uma função crescente em
$[0,\infty)$, obtemos finalmente a desigualdade \eqref{eq:lyapunov}.
\end{proof}

A ordem das duas últimas subseções não é acidental. Para aplicar Jensen a
$Y=Z^s$, precisamos saber que $\mathbb{E}[Z^s]$ é finita, e o
\Cref{lem:potencias} fornece essa informação a partir de $\mathbb{E}[Z^r]$ sem
usar Lyapunov. Se tentássemos provar a finitude por Lyapunov, usaríamos a
desigualdade que queremos demonstrar.

\begin{corollary}\label{cor:lyapunov-momentos}
Seja $X$ uma variável aleatória e sejam $0<s\leqslant r$. Então
\[
\bigl(\mathbb{E}|X|^s\bigr)^{1/s}
\leqslant
\bigl(\mathbb{E}|X|^r\bigr)^{1/r}.
\]
Em particular, $\mathbb{E}|X|\leqslant\bigl(\mathbb{E}[X^2]\bigr)^{1/2}$.
\end{corollary}

\begin{proof}
Aplique o \Cref{thm:lyapunov} a $Z=|X|$.
\end{proof}

A desigualdade de Lyapunov permite comparar as cotas obtidas aplicando a
desigualdade de Markov a diferentes potências de $|X|$. Suponha que
$\mathbb{E}|X|^r<\infty$, e escreva
$\|X\|_q=(\mathbb{E}|X|^q)^{1/q}$ para $0<q\leqslant r$. Para cada
$s\in(0,r]$ e $t>0$, a desigualdade de Markov dá
\[
\mathbb{P}(|X|\geqslant t)
=
\mathbb{P}(|X|^s\geqslant t^s)
\leqslant
\frac{\mathbb{E}|X|^s}{t^s}
=\left(\frac{\|X\|_s}{t}\right)^s.
\]
Pelo \Cref{cor:lyapunov-momentos}, $\|X\|_s\leqslant\|X\|_r$; logo, a
cota acima também é limitada por $(\|X\|_r/t)^s$. Quando
$t>\|X\|_r$, a razão $\|X\|_r/t$ está entre $0$ e $1$, de modo que, para
$s<r$,
\[
\left(\frac{\|X\|_r}{t}\right)^r
<
\left(\frac{\|X\|_r}{t}\right)^s.
\]
Assim, entre essas cotas comparadas usando a mesma quantidade $\|X\|_r$, a
que usa o maior expoente fornece o menor limite para desvios acima de
$\|X\|_r$. Em particular, usando diretamente o momento de ordem $r$,
obtemos
\[
\mathbb{P}(|X|\geqslant t)
\leqslant \frac{\mathbb{E}|X|^r}{t^r},
\]
uma cota que decai como $t^{-r}$ quando $t\to\infty$. Portanto, quanto maior
for a ordem do momento finito disponível, mais rápido pode ser o decaimento
polinomial garantido para a cauda de $|X|$.

\subsection{A Curva das Normas}\label{subsec:curva-normas}

O \Cref{thm:lyapunov} diz que a função $r\longmapsto(\mathbb{E}|X|^r)^{1/r}$ é não
decrescente. Há uma segunda estrutura por trás dela: o logaritmo dessa função é
convexo. A prova usa a desigualdade de Hölder, que é a forma geral da
desigualdade de Cauchy-Schwarz e que também sai da convexidade.

\begin{lemma}[Desigualdade de Young]\label{lem:young}
Sejam $p,q>1$ com $1/p+1/q=1$ e sejam $a,b\geqslant0$. Então
\[
ab\leqslant\frac{a^p}{p}+\frac{b^q}{q}.
\]
\end{lemma}

\begin{proof}
Se $a=0$ ou $b=0$, a desigualdade é imediata. Caso contrário, a função
$t\longmapsto e^t$ é convexa, logo
\[
ab=\exp\Bigl(\frac{1}{p}\log a^p+\frac{1}{q}\log b^q\Bigr)
\leqslant
\frac{1}{p}e^{\log a^p}+\frac{1}{q}e^{\log b^q}
=
\frac{a^p}{p}+\frac{b^q}{q}.
\]
\end{proof}

\begin{theorem}[Desigualdade de Hölder]\label{thm:holder}
Sejam $p,q>1$ com $1/p+1/q=1$ e sejam $X,Y$ variáveis aleatórias tais que
$\mathbb{E}|X|^p<\infty$ e $\mathbb{E}|Y|^q<\infty$. Então
\[
\mathbb{E}|XY|\leqslant\bigl(\mathbb{E}|X|^p\bigr)^{1/p}
\bigl(\mathbb{E}|Y|^q\bigr)^{1/q}.
\]
\end{theorem}

\begin{proof}
Se $\mathbb{E}|X|^p=0$, então $X=0$ quase certamente e a desigualdade é
imediata. O mesmo vale para $Y$. Suponha então que os dois momentos sejam
positivos e defina
\[
\tilde X=\frac{X}{(\mathbb{E}|X|^p)^{1/p}},
\qquad
\tilde Y=\frac{Y}{(\mathbb{E}|Y|^q)^{1/q}}.
\]
Então $\mathbb{E}|\tilde X|^p=\mathbb{E}|\tilde Y|^q=1$. Aplicando o
\Cref{lem:young} a $|\tilde X|$ e $|\tilde Y|$, obtemos
\[
|\tilde X\tilde Y|\leqslant\frac{|\tilde X|^p}{p}+\frac{|\tilde Y|^q}{q}.
\]
Tomando esperanças,
\[
\mathbb{E}|\tilde X\tilde Y|\leqslant\frac{1}{p}+\frac{1}{q}=1,
\]
que é a desigualdade para $X$ e $Y$.
\end{proof}

\begin{proposition}\label{prop:log-convexa}
Seja $X$ uma variável aleatória com $\mathbb{P}(X\neq0)>0$ e seja
$I=\{r>0:\mathbb{E}|X|^r<\infty\}$. Então a função $r\longmapsto\log\mathbb{E}|X|^r$
é convexa em $I$.
\end{proposition}

\begin{proof}
Sejam $r_1<r_2$ em $I$ e $t\in(0,1)$. Escreva
$r=(1-t)r_1+tr_2$. Queremos mostrar que
\[
\log\mathbb{E}|X|^r
\leqslant
(1-t)\log\mathbb{E}|X|^{r_1}+t\log\mathbb{E}|X|^{r_2}.
\]
Equivalentemente, depois de exponenciar,
\[
\mathbb{E}|X|^r
\leqslant
\bigl(\mathbb{E}|X|^{r_1}\bigr)^{1-t}
\bigl(\mathbb{E}|X|^{r_2}\bigr)^{t}.
\]
Escrevemos $|X|^r=|X|^{(1-t)r_1}|X|^{tr_2}$ e aplicamos a desigualdade de
Hölder com
\[
p=\frac{1}{1-t},
\qquad
q=\frac{1}{t},
\]
de modo que $1/p+1/q=1$. Tomando
$X_1=|X|^{(1-t)r_1}$ e $X_2=|X|^{tr_2}$, Hölder dá
\[
\mathbb{E}|X|^r
\leqslant
\bigl(\mathbb{E}|X|^{r_1}\bigr)^{1-t}
\bigl(\mathbb{E}|X|^{r_2}\bigr)^{t},
\]
como queríamos.
\end{proof}

\subsection{Aplicação: O Jogo Honesto e a Ruína}\label{subsec:jogo}

Considere o seguinte jogo. O jogador começa com capital $W_0=1$. Em cada
rodada, o capital é multiplicado por $3/2$ ou por $1/2$, com probabilidade
$1/2$ cada, independentemente das rodadas anteriores. O jogo é honesto no
sentido de que a taxa de retorno esperado, em cada uma das rodadas é dada por cópias independentes da variável 
aleatória $R_1$ cuja a distribuição é dada por 
\begin{align*}
\mathbb{P}\left( R_1=\frac{3}{2} \right)=\frac{1}{2} = \mathbb{P}\left( R_1=\frac{1}{2} \right)
\end{align*} 
e seu valor esperado é dado por
\[
\mathbb{E}[R_1]=\frac{1}{2}\cdot\frac{3}{2}+\frac{1}{2}\cdot\frac{1}{2}=1.
\]

Pela hipótese de independência, sabemos que taxa de retorno esperada deste jogo após $n$ rodadas
é dada pela esperança da variável aleatória $W_n\equiv R_1\cdot R_2\cdot\ldots\cdot R_n$ que é 
precisamente
\begin{align*}
\mathbb{E}[W_n]
=  
\mathbb{E}[R_1\cdot R_2\cdot\ldots\cdot R_n] 
=   
\mathbb{E}[R_1] \cdot \mathbb{E}[R_2]\cdot\ldots\cdot \mathbb{E}[R_n] 
=
1.
\end{align*}

A pergunta é se o
capital típico se mantém, cresce ou decai.
Para entender melhor este problema observe que o capital (aqui o capital inicial é de uma unidade) 
depois de $n$ rodadas é dado por 
\[
W_n=R_1R_2\cdots R_n,
\]
com os $R_i$ independentes e identicamente distribuídos. Como $W_n>0$ sempre,
podemos tomar logaritmos. A convexidade de $-\log$ ou, equivalentemente, a
concavidade de $\log$, dá
\[
\mathbb{E}[\log W_n]\leqslant\log\mathbb{E}[W_n]=0.
\]
A desigualdade é estrita, porque $W_n$ não é constante. De fato,
\[
\mathbb{E}[\log R_1]
=
\frac{1}{2}\log\frac{3}{2}+\frac{1}{2}\log\frac{1}{2}
=
\frac{1}{2}\log\frac{3}{4}
=
-\frac{1}{2}\log\frac{4}{3}
\approx-0{,}14384,
\]
logo
\begin{equation}\label{eq:jogo-log}
\mathbb{E}[\log W_n]
=
\sum_{i=1}^{n}\mathbb{E}[\log R_i]
=
-\frac{n}{2}\log\frac{4}{3}
\;\xrightarrow[n\to\infty]{}\;-\infty.
\end{equation}
Enquanto o logaritmo da esperança do capital é zero, a esperança do logaritmo
tende a $-\infty$. A diferença entre as duas quantidades se manifesta como
sendo a diferença dos dois termos que aparecem na Desigualdade de Jensen. 
Neste caso vimos que esta diferença cresce linearmente em $n$.

Vale registrar a leitura pela desigualdade de Lyapunov. Como os $R_i$ são
independentes,
\[
\mathbb{E}[W_n^2]=\prod_{i=1}^{n}\mathbb{E}[R_i^2]
=\Bigl(\frac{5}{4}\Bigr)^{n},
\]
e o \Cref{cor:lyapunov-momentos} dá
\[
\mathbb{E}[W_n]\leqslant\bigl(\mathbb{E}[W_n^2]\bigr)^{1/2}
=\Bigl(\frac{5}{4}\Bigr)^{n/2}.
\]
A cota cresce exponencialmente, enquanto a esperança é constante. A distância entre
os dois valores médios apresenta também uma certa assimetria  
analogamente ao que vimos pela aplicação da Desigualdade de Jensen.

Para transformar \eqref{eq:jogo-log} em uma afirmação sobre o capital, sem
usar a lei dos grandes números, basta conhecer a distribuição exata. Seja
$K_n$ o número de rodadas em que o capital foi multiplicado por $3/2$. Então
$K_n\sim\mathrm{Bin}(n,1/2)$ e
\[
W_n=\Bigl(\frac{3}{2}\Bigr)^{K_n}\Bigl(\frac{1}{2}\Bigr)^{n-K_n}
=
\frac{3^{K_n}}{2^n}.
\]
Se $K_n\leqslant n/2$, então
\[
W_n\leqslant\frac{3^{n/2}}{2^n}
=
\Bigl(\frac{\sqrt3}{2}\Bigr)^{n}.
\]
Como a distribuição de $K_n$ é simétrica em torno de $n/2$, o evento
$\{K_n\leqslant n/2\}$ tem probabilidade pelo menos $1/2$. 
% Bin(11,1/2): alturas normalizadas pelo máximo 462/2048.
\begin{figure}[H]
\centering
\begin{tikzpicture}[x=0.85cm, y=3.1cm, font=\small, >=latex]
  % Duas metades espelhadas, separadas por uma abscissa não inteira.
  \foreach \k/\h in {0/0.002165,1/0.023810,2/0.119048,
                     3/0.357143,4/0.714286,5/1} {
    \filldraw[fill=teal!65, draw=teal, rounded corners=1pt]
      ({\k-0.32},0) rectangle ({\k+0.32},\h);
    \pgfmathsetmacro{\espelho}{11-\k}
    \filldraw[fill=muted!20, draw=muted!65, rounded corners=1pt]
      ({\espelho-0.32},0) rectangle ({\espelho+0.32},\h);
  }

  \draw[thin, ->, draw=muted] (-0.65,0) -- (11.85,0)
    node[right] {$k$};
  \foreach \k in {0,...,11} {
    \draw[draw=muted] (\k,0) -- (\k,-0.025);
    \node[below] at (\k,-0.025) {$\k$};
  }
  \draw[densely dashed, draw=muted!70] (5.5,0) -- (5.5,1.18);
  \node[above, text=muted] at (5.5,1.18) {$\frac{n}{2}=5{,}5$};

  % Cada par de barras simétricas tem a mesma massa.
  \draw[decorate, decoration={brace, mirror, amplitude=5pt}, draw=teal]
    (-0.32,-0.19) -- (5.32,-0.19)
    node[midway, below=8pt, text=teal]
      {$\mathbb{P}(K_{11}\leqslant5)=\dfrac12$};
  \draw[decorate, decoration={brace, mirror, amplitude=5pt}, draw=muted]
    (5.68,-0.19) -- (11.32,-0.19)
    node[midway, below=8pt, text=muted]
      {$\mathbb{P}(K_{11}\geqslant6)=\dfrac12$};
\end{tikzpicture}
\caption{Função de probabilidade de uma v.a com distribuição $\mathrm{Bin}(11,\tfrac{1}{2})$}
\label{fig:jogo-binomial-simetria}
\end{figure}

Logo
\begin{equation}\label{eq:jogo-mediana}
\mathbb{P}\Bigl(W_n\leqslant\Bigl(\frac{\sqrt3}{2}\Bigr)^{n}\Bigr)
\geqslant\frac{1}{2}.
\end{equation}
Com probabilidade pelo menos um meio, o capital decai a uma taxa geométrica de
$\sqrt3/2\approx0{,}866$ por rodada, isto é, pelo menos $13{,}4\%$ ao longo de
cada rodada, em termos geométricos. Quando $n$ é par, o valor mais provável de
$W_n$ ocorre em $K_n=n/2$ e vale exatamente $(\sqrt3/2)^n$. Quando $n$ é
ímpar, os dois valores mais prováveis correspondem a
$K_n=(n-1)/2$ e $K_n=(n+1)/2$.

A desigualdade de Chebyshev dá uma segunda leitura, agora sobre a
probabilidade de o jogador terminar acima do capital inicial. O evento
$\{W_n>1\}$ é o mesmo que $\{3^{K_n}>2^n\}$, ou seja,
\[
K_n>n\frac{\log2}{\log3}.
\]
Como $\log2/\log3\approx0{,}63093$, o evento é
$\{K_n-n/2>0{,}13093\,n\}$. A variância de $K_n$ é $n/4$, logo Chebyshev dá
\begin{equation}\label{eq:jogo-chebyshev}
\mathbb{P}(W_n>1)
\leqslant
\frac{n/4}{(0{,}13093\,n)^2}
=
\frac{14{,}58}{n}.
\end{equation}
Assim a probabilidade de terminar com mais do que se começou tende a zero,
embora a esperança do capital seja $1$ para todo $n$.

Os números para $n=20$ resumem o fenômeno. Temos $\mathbb{E}[W_{20}]=1$ e
$\mathbb{E}[W_{20}^2]=(5/4)^{20}\approx86{,}7$, de modo que o desvio padrão de
$W_{20}$ é aproximadamente $9{,}3$. A cota \eqref{eq:jogo-mediana} diz que, com
probabilidade pelo menos $1/2$, o capital é no máximo
$(0{,}866)^{20}\approx0{,}056$. O valor mais provável é exatamente esse, com
probabilidade $\binom{20}{10}2^{-20}\approx0{,}176$. Em outras palavras, o
capital mais provável é cerca de $5{,}6\%$ do inicial, enquanto a média é
$100\%$ do inicial. A média é sustentada por uma minoria de trajetórias com
ganhos muito grandes.

Como os fatores $R_1,\ldots,R_n$ são independentes, a esperança do produto é
o produto das esperanças e, neste jogo, $\mathbb{E}[W_n]=1$. Essa quantidade
descreve o capital médio, mas não basta para descrever como o capital evolui
na maioria das trajetórias. Para estudar o crescimento multiplicativo típico
é mais eficiente olhar para logaritmos  $\log W_n=\sum_{j=1}^n\log R_j$. Assim,
$\mathbb{E}[\log R_1]$ é o incremento médio por rodada na escala logarítmica;
como esse valor é negativo, os cálculos com a distribuição de $K_n$ acima
mostram que o capital costuma decair, embora sua esperança permaneça igual a
$1$. A desigualdade de Jensen expressa a diferença entre essas duas
quantidades:
\[
\mathbb{E}[\log W_n]<\log\mathbb{E}[W_n]=0.
\]
Portanto, a esperança do capital e o crescimento na escala logarítmica dão
informações diferentes: a esperança do capital é sustentada pelas trajetórias
raras em que ele cresce muito, enquanto o crescimento na escala logarítmica
registra o efeito médio dos fatores multiplicativos e, neste jogo, reflete o
decaimento típico do capital.

\section{O Lema de Borel-Cantelli}\label{sec:borel-cantelli}

Até aqui, estimamos a probabilidade de um evento. Nesta seção, a pergunta é
diferente: dada uma sequência de eventos $A_1,A_2,\ldots$, queremos saber se
eles ocorrem infinitas vezes. A soma $\sum_n\mathbb{P}(A_n)$ é o número
esperado de ocorrências, e o lema de Borel-Cantelli diz que o comportamento
dessa série decide a questão, desde que os eventos sejam independentes na
direção difícil.

\subsection{Limites Superior e Inferior de Eventos}\label{subsec:limsup}

Para formular a pergunta, precisamos de uma noção de limite para sequências de
eventos.

\begin{definition}[Limite superior]\label{def:limsup}
Seja $(A_n)_{n\geqslant1}$ uma sequência de eventos em um espaço de
probabilidade $(\Omega,\mathcal{F},\mathbb{P})$. O \emph{limite superior} de
$(A_n)$ é o evento
\[
\limsup_{n\to\infty}A_n
=
\bigcap_{n=1}^{\infty}\bigcup_{k=n}^{\infty}A_k.
\]
Escrevemos também $\{A_n\text{ infinitas vezes}\}$ ou, abreviadamente,
$\{A_n\ \text{i.v.}\}$ para esse evento.
\end{definition}

A definição merece uma leitura em palavras. Um ponto $\omega$ pertence a
$\bigcup_{k\geqslant n}A_k$ se, e somente se, $\omega$ pertence a algum $A_k$
com $k\geqslant n$. Se $\omega$ pertence à interseção sobre $n$, então, para
cada $n$, existe $k\geqslant n$ com $\omega\in A_k$. Isso significa que os
índices $k$ para os quais $\omega\in A_k$ não são limitados, ou seja, $\omega$
pertence a infinitos dos eventos $A_n$. Reciprocamente, se $\omega$ pertence a
infinitos $A_n$, então para cada $n$ existe $k\geqslant n$ com
$\omega\in A_k$, e $\omega$ está no limite superior. Assim
\[
\limsup_{n\to\infty}A_n
=
\{\omega\in\Omega:\omega\in A_n\text{ para infinitos índices }n\}.
\]

O limite inferior tem definição dual.

\begin{definition}[Limite inferior]\label{def:liminf}
Nas mesmas hipóteses, o \emph{limite inferior} de $(A_n)$ é o evento
\[
\liminf_{n\to\infty}A_n
=
\bigcup_{n=1}^{\infty}\bigcap_{k=n}^{\infty}A_k.
\]
Escrevemos também $\{A_n\text{ eventualmente}\}$.
\end{definition}

Aqui, $\omega$ pertence a $\bigcap_{k\geqslant n}A_k$ se, e somente se,
$\omega\in A_k$ para todo $k\geqslant n$. Se $\omega$ pertence à união sobre
$n$, então existe $n$ tal que $\omega\in A_k$ para todo $k\geqslant n$, isto é,
$\omega$ pertence a todos os $A_n$ a partir de algum índice. Logo
\[
\liminf_{n\to\infty}A_n
=
\{\omega\in\Omega:\omega\in A_n\text{ para todo }n\text{ suficientemente grande}\}.
\]

Os dois eventos se relacionam por complementação. Como o complementar troca
uniões por interseções e inverte o sentido das inclusões,
\begin{equation}\label{eq:limsup-complementar}
\Bigl(\limsup_{n\to\infty}A_n\Bigr)^c
=
\liminf_{n\to\infty}A_n^c.
\end{equation}
Além disso, todo ponto que pertence a todos os $A_n$ a partir de um índice
pertence a infinitos deles, logo
\[
\liminf_{n\to\infty}A_n
\subseteq
\limsup_{n\to\infty}A_n.
\]
A inclusão pode ser estrita: no lançamento de uma moeda, o evento ``sai cara
infinitas vezes'' contém o evento ``sai cara a partir de algum lançamento'', e a
inclusão é estrita quase certamente.

\begin{remark}\label{rem:limsup-indicadoras}
A definição se traduz em uma identidade entre indicadoras:
\[
\mathds{1}_{\limsup_n A_n}
=
\limsup_{n\to\infty}\mathds{1}_{A_n},
\qquad
\mathds{1}_{\liminf_n A_n}
=
\liminf_{n\to\infty}\mathds{1}_{A_n}.
\]
Em particular, $\omega\in\limsup_n A_n$ se, e somente se,
\[
\sum_{n=1}^{\infty}\mathds{1}_{A_n}(\omega)=\infty.
\]
Essa leitura é útil porque transforma a pergunta sobre eventos em uma pergunta
sobre séries de números não-negativos.
\end{remark}

\subsection{A Continuidade das Medidas de Probabilidade}\label{subsec:continuidade}

O instrumento técnico que liga os limites de eventos às probabilidades é a
continuidade das medidas, que já foi demonstrada na nota de propriedades de
probabilidade. Recordamos o enunciado.

\begin{theorem}[Continuidade para sequências crescentes]\label{thm:continuidade-crescente}
Seja $(\Omega,\mathcal{F},\mathbb{P})$ um espaço de probabilidade e seja
$(A_n)_{n\geqslant1}$ uma sequência crescente de eventos, com
$A=\bigcup_{n\geqslant1}A_n$. Então
\[
\mathbb{P}(A)=\lim_{n\to\infty}\mathbb{P}(A_n).
\]
\end{theorem}

\begin{corollary}[Continuidade para sequências decrescentes]\label{cor:continuidade-decrescente}
Seja $(B_n)_{n\geqslant1}$ uma sequência decrescente de eventos, com
$B=\bigcap_{n\geqslant1}B_n$. Então
\[
\mathbb{P}(B)=\lim_{n\to\infty}\mathbb{P}(B_n).
\]
\end{corollary}

A prova do teorema está na nota de propriedades de probabilidade, onde a união
é decomposta em anéis disjuntos. O corolário segue aplicando o teorema aos
complementares, que formam uma sequência crescente de limite
$\Omega\setminus B$.

No que segue, aplicaremos o corolário à sequência
\[
C_n=\bigcup_{k=n}^{\infty}A_k.
\]
Essa sequência é decrescente, pois $C_{n+1}\subseteq C_n$, e sua interseção é
exatamente $\limsup_n A_n$. Logo
\begin{equation}\label{eq:limsup-continuidade}
\mathbb{P}\Bigl(\limsup_{n\to\infty}A_n\Bigr)
=
\lim_{n\to\infty}\mathbb{P}(C_n).
\end{equation}
A identidade \eqref{eq:limsup-continuidade} é a ponte entre o lema de
Borel-Cantelli e a continuidade.

\subsection{A Desigualdade de Boole}\label{subsec:boole}

Falta controlar a probabilidade de uma união enumerável. A desigualdade de
Boole faz isso, e sua versão finita já apareceu na nota de propriedades de
probabilidade.

\begin{lemma}[Desigualdade de Boole]\label{lem:boole}
Seja $(A_n)_{n\geqslant1}$ uma sequência de eventos. Então
\[
\mathbb{P}\Bigl(\bigcup_{n=1}^{\infty}A_n\Bigr)
\leqslant
\sum_{n=1}^{\infty}\mathbb{P}(A_n).
\]
\end{lemma}

\begin{proof}
Defina $B_1=A_1$ e, para $n\geqslant2$,
\[
B_n=A_n\setminus\bigcup_{k=1}^{n-1}A_k.
\]
Os eventos $B_n$ são dois a dois disjuntos e
$\bigcup_{n\geqslant1}B_n=\bigcup_{n\geqslant1}A_n$. Além disso,
$B_n\subseteq A_n$, logo $\mathbb{P}(B_n)\leqslant\mathbb{P}(A_n)$. Pela
$\sigma$-aditividade,
\[
\mathbb{P}\Bigl(\bigcup_{n=1}^{\infty}A_n\Bigr)
=
\sum_{n=1}^{\infty}\mathbb{P}(B_n)
\leqslant
\sum_{n=1}^{\infty}\mathbb{P}(A_n).
\]
\end{proof}

\subsection{O Primeiro Lema de Borel-Cantelli}\label{subsec:bc1}

Vermos a seguir que a continuidade da medida e a desigualdade de Boole
é praticamente tudo que precisamos para a prova da primeira parte do Lema
de Borel-Cantelli.

\begin{theorem}[Primeiro lema de Borel-Cantelli]\label{thm:bc1}
Seja $(A_n)_{n\geqslant1}$ uma sequência de eventos tal que
$\sum_{n\geqslant1}\mathbb{P}(A_n)<\infty$. Então
\[
\mathbb{P}\Bigl(\limsup_{n\to\infty}A_n\Bigr)=0.
\]
Em outras palavras, com probabilidade $1$ ocorrem apenas finitos dos eventos
$A_n$.
\end{theorem}

\begin{proof}
Já que a sequência de eventos $C_n\equiv \bigcup_{k\geqslant n}A_k$
é uma sequência decrescente de eventos, segue da continuidade da medida que 
\[
\mathbb{P}\Bigl(\limsup_{n\to\infty}A_n\Bigr)
=
\lim_{n\to\infty}\mathbb{P}(C_n).
\]
Portanto segue da Desigualdade de Boole \Cref{lem:boole} que
\[
\mathbb{P}(C_n)
\leqslant
\sum_{k=n}^{\infty}\mathbb{P}(A_k).
\]
Como a série $\sum_k\mathbb{P}(A_k)$ converge, sua cauda
$\sum_{k\geqslant n}\mathbb{P}(A_k)\to 0$, quando $n\to\infty$. Logo
$\mathbb{P}(C_n)\to0$ e, portanto,
$\mathbb{P}(\limsup_n A_n)=0$.
\end{proof}

Note que a prova não usou independência. O primeiro lema vale para qualquer
sequência de eventos. 

\subsection{O Segundo Lema de Borel-Cantelli}\label{subsec:bc2}

A recíproca do primeiro lema é falsa em geral: se $A_n=A$ para todo $n$, com
$0<\mathbb{P}(A)<1$, então $\sum_n\mathbb{P}(A_n)=\infty$ e
$\limsup_n A_n=A$, cuja probabilidade não é $1$. O segundo lema mostra que a
recíproca vale quando os eventos são independentes.

Usaremos a desigualdade elementar
\begin{equation}\label{eq:exp-ineq}
1-x\leqslant e^{-x}
\qquad
\text{para }x\in\mathbb{R},
\end{equation}
que segue da convexidade de $x\longmapsto e^x$, ou da comparação das séries de
Taylor.

\begin{theorem}[Segundo lema de Borel-Cantelli]\label{thm:bc2}
Seja $(A_n)_{n\geqslant1}$ uma sequência de eventos independentes tal que
$\sum_{n\geqslant1}\mathbb{P}(A_n)=\infty$. Então
\[
\mathbb{P}\Bigl(\limsup_{n\to\infty}A_n\Bigr)=1.
\]
\end{theorem}

\begin{proof}
Seja $p_n=\mathbb{P}(A_n)$. Pela identidade \eqref{eq:limsup-complementar},
\[
\Bigl(\limsup_{n\to\infty}A_n\Bigr)^c
=
\liminf_{n\to\infty}A_n^c
=
\bigcup_{m=1}^{\infty}\bigcap_{n=m}^{\infty}A_n^c.
\]
Pela desigualdade de Boole, basta mostrar que
$\mathbb{P}\bigl(\bigcap_{n\geqslant m}A_n^c\bigr)=0$ para cada $m$.

Fixe $m$. Para $N\geqslant m$, os eventos $A_m^c,\ldots,A_N^c$ são
independentes, logo
\[
\mathbb{P}\Bigl(\bigcap_{n=m}^{N}A_n^c\Bigr)
=
\prod_{n=m}^{N}(1-p_n)
\leqslant
\exp\Bigl(-\sum_{n=m}^{N}p_n\Bigr),
\]
onde usamos a desigualdade elementar \eqref{eq:exp-ineq} em cada fator. A sequência
$\bigcap_{n=m}^{N}A_n^c$ é decrescente em $N$, com interseção
$\bigcap_{n\geqslant m}A_n^c$. Pela propriedade de continuidade da medida
\Cref{cor:continuidade-decrescente}, podemos afirmar que
\[
\mathbb{P}\Bigl(\bigcap_{n=m}^{\infty}A_n^c\Bigr)
=
\lim_{N\to\infty}\mathbb{P}\Bigl(\bigcap_{n=m}^{N}A_n^c\Bigr)
\leqslant
\lim_{N\to\infty}\exp\Bigl(-\sum_{n=m}^{N}p_n\Bigr)
=
0,
\]
onde a última igualdade segue da hipótese da série $\sum_n p_n$ ser divergente. 
Como $m$ é arbitrário, todos os eventos
$\bigcap_{n\geqslant m}A_n^c$ têm probabilidade nula, e a união enumerável
deles também. Logo $\mathbb{P}\bigl((\limsup_n A_n)^c\bigr)=0$, isto é,
$\mathbb{P}(\limsup_n A_n)=1$.
\end{proof}

\begin{remark}\label{rem:bc2-independencia}
A hipótese de independência é indispensável no segundo lema. O exemplo
$A_n=A$, com $0<\mathbb{P}(A)<1$, mostra que sem ela a conclusão falha. Um
segundo exemplo, que mostra que a conclusão pode falhar mesmo quando os eventos
são distintos, é $A_n=[0,1/n]$ no espaço $([0,1],\mathcal{B}([0,1]),\lambda)$,
onde $\lambda$ é a medida uniforme: a soma das probabilidades diverge, mas a
interseção dos $A_n$ é $\{0\}$, de modo que $\limsup_n A_n=\{0\}$ e tem
probabilidade nula.
\end{remark}

\subsection{A Lei 0-1 de Borel}\label{subsec:lei-01}

Reunindo as duas partes, obtemos um critério completo para eventos
independentes.

\begin{corollary}[Lei 0-1 de Borel]\label{cor:lei-01}
Seja $(A_n)_{n\geqslant1}$ uma sequência de eventos independentes. Então
\[
\mathbb{P}\Bigl(\limsup_{n\to\infty}A_n\Bigr)\in\{0,1\}.
\]
Mais precisamente, a probabilidade é $0$ se $\sum_n\mathbb{P}(A_n)<\infty$ e
igual a $1$ se $\sum_n\mathbb{P}(A_n)=\infty$.
\end{corollary}

\begin{proof}
Se a série converge, o \Cref{thm:bc1} dá probabilidade $0$. Se diverge, o
\Cref{thm:bc2} dá probabilidade $1$.
\end{proof}

\subsection{Exemplos}\label{subsec:bc-exemplos}

\begin{example}[Caras em lançamentos sucessivos]\label{ex:caras-infinitas}
Lance uma moeda honesta infinitas vezes e seja $A_n$ o evento ``sai cara no
lançamento $n$''. Os eventos são independentes, $\mathbb{P}(A_n)=1/2$ e
$\sum_n\mathbb{P}(A_n)=\infty$. Pelo \Cref{thm:bc2}, saem caras infinitas vezes
com probabilidade $1$.
\end{example}

\begin{example}[Blocos iniciais de caras]\label{ex:blocos-iniciais}
No mesmo espaço, seja $B_n$ o evento ``os primeiros $n$ lançamentos são todos
caras''. Então $\mathbb{P}(B_n)=2^{-n}$ e $\sum_n2^{-n}<\infty$. Pelo
\Cref{thm:bc1}, com probabilidade $1$ apenas finitos dos $B_n$ ocorrem. Isso
\textbf{não} contradiz o exemplo anterior: o que se repete infinitas vezes são caras em
posições isoladas, não blocos de caras que começam na posição $1$.
\end{example}



\subsection{Aplicação: A Maior Sequência de Caras}\label{subsec:maior-sequencia}

Considere uma moeda honesta lançada infinitas vezes e seja $L_n$ o comprimento
da maior sequência de caras consecutivas inteiramente contida nos primeiros $n$
lançamentos. Quanto vale $L_n$ para $n$ grande? A intuição sugere que
sequências longas da mesma face são raras, entretanto o número de posições onde elas podem começar
cresce com $n$, e o equilíbrio entre os dois efeitos produz a escala
$\log_2 n$. O lema de Borel-Cantelli determina a constante.

Seja $A_i$ o evento ``sai cara no lançamento $i$''. Para $k\geqslant1$ e
$j\geqslant1$, o evento
\[
B_{j,k}=A_j\cap A_{j+1}\cap\cdots\cap A_{j+k-1}
\]
tem probabilidade $2^{-k}$, e eventos desse tipo são independentes quando os
blocos de índices são disjuntos.

\smallskip
\noindent\textbf{Limite superior.}
Fixe $\varepsilon>0$ e, para cada $m\geqslant2$, ponha
$k_m=\lceil(1+\varepsilon)\log_2 m\rceil$. Para $m\geqslant k_m$, seja
$E_m=B_{m-k_m+1,k_m}$ o evento de que uma sequência de $k_m$ caras consecutivas
termine exatamente no lançamento $m$. Então
\[
\mathbb{P}(E_m)
=
2^{-k_m}
\leqslant
m^{-(1+\varepsilon)}.
\]
Como $1+\varepsilon>1$, a série $\sum_{m}m^{-(1+\varepsilon)}$ converge, e o
\Cref{thm:bc1} garante que, quase certamente, apenas finitos dos eventos $E_m$
ocorrem. Existe, portanto, um índice aleatório $N_0$ tal que $E_m$ não ocorre
para nenhum $m\geqslant N_0$.

Considere agora $n\geqslant N_0$ e qualquer sequência de caras consecutivas
contida em $\{1,\ldots,n\}$, digamos terminando na posição $m\leqslant n$. Se
$m<N_0$, seu comprimento é no máximo $m<N_0$; se $N_0\leqslant m\leqslant n$,
como $E_m$ não ocorre, seu comprimento é estritamente menor do que
$k_m\leqslant k_n$. Logo $L_n\leqslant\max\{N_0,k_n-1\}$. Para $n$
suficientemente grande temos $k_n>N_0$, donde
$L_n\leqslant k_n-1<(1+\varepsilon)\log_2 n$. Em particular,
\[
\limsup_{n\to\infty}\frac{L_n}{\log_2 n}\leqslant1
\qquad
\text{quase certamente}.
\]

\smallskip
\noindent\textbf{Limite inferior.}
Fixe $\varepsilon\in(0,1)$ e ponha
$k_n=\lfloor(1-\varepsilon)\log_2 n\rfloor$. Dividimos os primeiros $n$
lançamentos em $m_n=\lfloor n/k_n\rfloor$ blocos disjuntos de $k_n$ posições
consecutivas. Se $L_n<k_n$, nenhum desses blocos é todo de caras. Como os
blocos são disjuntos, os eventos correspondentes são independentes, e cada um
tem probabilidade $2^{-k_n}$ de ser todo de caras. Logo
\[
\mathbb{P}(L_n<k_n)
\leqslant
\bigl(1-2^{-k_n}\bigr)^{m_n}
\leqslant
\exp\bigl(-m_n2^{-k_n}\bigr).
\]
Agora $2^{-k_n}\geqslant n^{-(1-\varepsilon)}$ e, para $n$ grande,
$m_n\geqslant n/(2k_n)$. Assim
\[
m_n2^{-k_n}
\geqslant
\frac{n^{\varepsilon}}{2k_n}
\geqslant
\frac{n^{\varepsilon}}{2(1+\log_2 n)}
\;\xrightarrow[n\to\infty]{}\;\infty.
\]
Segue que $\mathbb{P}(L_n<k_n)\leqslant\exp\bigl(-n^{\varepsilon}/(2k_n)\bigr)$,
e a série dessas probabilidades converge, porque a exponencial decai mais
rápido do que qualquer potência de $n$. Pelo \Cref{thm:bc1},
\[
\mathbb{P}\bigl(L_n<(1-\varepsilon)\log_2 n
\text{ para infinitos }n\bigr)=0,
\]
isto é, quase certamente $L_n\geqslant k_n$ para todo $n$ suficientemente
grande. Logo
\[
\liminf_{n\to\infty}\frac{L_n}{\log_2 n}\geqslant1
\qquad
\text{quase certamente}.
\]

\smallskip
\noindent\textbf{Conclusão.}
Reunindo os dois limites,
\begin{equation}\label{eq:maior-sequencia}
\frac{L_n}{\log_2 n}\;\xrightarrow[n\to\infty]{}\;1
\qquad
\text{quase certamente}.
\end{equation}
Mais precisamente, para cada $\varepsilon>0$ e quase todo $\omega$, vale
\[
(1-\varepsilon)\log_2 n-1\leqslant L_n\leqslant(1+\varepsilon)\log_2 n
\]
para todo $n$ suficientemente grande.

Vale observar o papel das duas direções. O limite superior controla a
probabilidade de uma sequência longa terminar no instante $m$ e usa o primeiro
lema para concluir que isso só ocorre um número finito de vezes. O limite
inferior usa a independência de blocos disjuntos para estimar a probabilidade de
não haver nenhum bloco de caras, e depois aplica novamente o primeiro lema. A
segunda parte do lema não é usada, o que mostra que ela não é a única forma de
obter conclusões de probabilidade $1$.

Os números tornam a escala concreta. Como $\log_2(10^6)\approx19{,}93$, em um
milhão de lançamentos a maior sequência tem cerca de $20$ caras. Para ver $30$
caras seguidas é preciso da ordem de $2^{30}\approx10^9$ lançamentos, e para
$40$ caras, da ordem de $2^{40}\approx10^{12}$. A probabilidade de $20$ caras
em $20$ lançamentos é $2^{-20}\approx9{,}5\times10^{-7}$, mas há cerca de
$10^6$ posições onde a sequência pode começar, e o produto é da ordem de $1$.

%----------------------------------------------------------------------
\section{Exercícios Propostos}\label{sec:exercicios}

\begin{exercise}\label{exercise:mediana-igualdade}
Determine quais distribuições em $[0,\infty)$, com média $\mu$, têm $2\mu$
como mediana. Isto é, determine quando a cota da
\Cref{prop:mediana-markov} é atingida.
\end{exercise}

\begin{exercise}\label{exercise:chebyshev-dado}
Use a desigualdade de Chebyshev para mostrar que, em $n$ lançamentos de um dado
honesto, a probabilidade de o número de seis estar entre
$\frac{1}{6}n-\sqrt n$ e $\frac{1}{6}n+\sqrt n$ é pelo menos $\frac{31}{36}$.
\end{exercise}

\begin{starredexercise}[Desigualdade de Cantelli]\label{exercise:chebyshev-forma}
Seja $X$ uma variável aleatória com média $\mu$ e variância $\sigma^2<\infty$.
Mostre que, para todo $a>0$,
\[
\mathbb{P}(X-\mu\geqslant a)\leqslant\frac{\sigma^2}{\sigma^2+a^2}.
\]
\emph{Sugestão:} Para $c>0$, observe que $X-\mu\geqslant a$ implica
$X-\mu+c\geqslant a+c$, aplique a desigualdade de Markov a $(X-\mu+c)^2$ e
minimize a cota em $c$.
\end{starredexercise}

\begin{exercise}\label{exercise:segunda-derivada}
Seja $I$ um intervalo aberto da reta e seja $f\colon I\to\mathbb{R}$ duas vezes
diferenciável, com $f''(x)>0$ para todo $x\in I$. Usando a fórmula de Taylor
com resto de Lagrange, mostre que $f$ é estritamente convexa em $I$. Dê um
exemplo mostrando que a recíproca não vale.
\end{exercise}

\begin{exercise}[Entropia de Shannon]\label{exercise:jensen-quadrado}
Sejam $p_1,\ldots,p_n>0$ tais que $\sum_{i=1}^{n}p_i=1$. Usando a desigualdade
de Jensen, mostre que
\[
-\sum_{i=1}^{n}p_i\log p_i\leqslant\log n,
\]
e determine quando vale a igualdade.
\end{exercise}

\begin{exercise}\label{exercise:media-harmonica}
Sejam $x_1,\ldots,x_n>0$ e seja $\eta$ a média harmônica definida no
\Cref{ex:media-harmonica}. Aplicando diretamente a desigualdade de Jensen à
função estritamente convexa $g(x)=1/x$ em $(0,\infty)$, mostre que
$\eta\leqslant\bar{x}$, onde $\bar{x}$ é a média aritmética, e determine quando
vale a igualdade.
\end{exercise}

\begin{exercise}\label{exercise:lyapunov-potencia}
Seja $p\geqslant1$. Usando a desigualdade de Jensen, encontre o menor valor da
constante $c_p$ tal que
\[
(a+b)^p\leqslant c_p(a^p+b^p)
\qquad
\text{para todos }a,b\geqslant0.
\]
\end{exercise}

\begin{exercise}\label{exercise:supremo-essencial}
Seja $X$ uma variável aleatória limitada e defina
$M=\inf\{c\geqslant0:\mathbb{P}(|X|\leqslant c)=1\}$. Mostre que
\[
\lim_{r\to\infty}\bigl(\mathbb{E}|X|^r\bigr)^{1/r}=M.
\]
\end{exercise}

\begin{exercise}\label{exercise:bc-produto}
Sejam $A_1,A_2,\ldots$ eventos independentes com $\mathbb{P}(A_n)=p_n\in[0,1)$
para todo $n$. Mostre que
\[
\mathbb{P}\Bigl(\bigcap_{n=1}^{\infty}A_n^c\Bigr)>0
\]
se, e somente se, $\sum_n p_n<\infty$.
\end{exercise}

\begin{starredexercise}\label{exercise:bc-dois-a-dois}
Mostre que a conclusão do \Cref{thm:bc2} continua válida se a hipótese de
independência for substituída por independência dois a dois. Isto é, mostre
que, se $(A_n)$ é uma sequência de eventos dois a dois independentes com
$\sum_n\mathbb{P}(A_n)=\infty$, então
$\mathbb{P}(\limsup_n A_n)=1$.
\emph{Sugestão:} Aplique a desigualdade de Chebyshev a
$S_{m,N}=\sum_{n=m}^{N}\mathds{1}_{A_n}$.
\end{starredexercise}

\begin{exercise}\label{exercise:bc-exponencial}
Sejam $X_1,X_2,\ldots$ variáveis aleatórias independentes com distribuição
exponencial de parâmetro $1$ e seja $M_n=\max\{X_1,\ldots,X_n\}$. Mostre que
\[
\frac{M_n}{\log n}\;\xrightarrow[n\to\infty]{}\;1
\qquad
\text{quase certamente}.
\]
\end{exercise}

%----------------------------------------------------------------------
% BIBLIOGRAFIA
%----------------------------------------------------------------------
\bibliographystyle{plain}
\bibliography{../arquivos-aux/referencias}

\end{document}
